% Generated human-review packet template.  Source records, Lean statements,
% and raw-source-to-Spec screening fields are substituted by
% scripts/review_dashboard_packet.py.
% Do not hand-edit generated packets; annotate the PDF or reconcile notes into
% the paper-local human-review log.
\documentclass[10pt]{article}
\usepackage[margin=0.43in,footskip=0.2in]{geometry}
\usepackage{fontspec}
% Use a Unicode-complete main face as well as a Unicode-complete code face.
% Reviewer reasons and source labels can contain exact Greek symbols outside
% verbatim blocks; silently dropping one would corrupt the review packet.
\setmainfont{DejaVu Serif}
% The broad DejaVu Sans Unicode coverage preserves Lean's mathematical glyphs
% (including U+2216 set difference and superscript identifiers) in verbatim
% review blocks.  Its proportional metrics are preferable to silently missing
% exact semantic text from a narrower monospace font.
\setmonofont{DejaVu Sans}
\usepackage{hyperref}
\usepackage{amssymb}
\usepackage{fvextra}
\usepackage{enumitem}
\setlength{\parindent}{0pt}
\setlength{\parskip}{0.15em}
\linespread{0.92}
\hypersetup{colorlinks=true,linkcolor=blue,urlcolor=blue,breaklinks=true}
\DefineVerbatimEnvironment{ReviewVerbatim}{Verbatim}{breaklines=true,breakanywhere=true,fontsize=\scriptsize}
\newcommand{\reviewerbox}{%
  \par\noindent\fbox{\parbox{0.96\linewidth}{\rule{0pt}{0.38in}}}\par
}
\newcommand{\reviewmatch}[1]{%
  \par\noindent\textbf{Reviewer check:} $\square$\; #1\par
  \noindent\emph{If left unchecked, explain the discrepancy or uncertainty below.}\par
}

\begin{document}
\fontsize{9}{10}\selectfont

\begin{center}
{\LARGE Human review packet: HT26EFXChores}\\[0.35em]
\small Generated from byte-pinned source excerpts, the expanded PaperInterface
specifications, and hash-bound raw-source-to-Spec screening records.
\end{center}

\noindent\textbf{Paper:} HT26EFXChores\\
\textbf{Paper id:} \texttt{HT26EFXChores}\\
\textbf{Source version:} \parbox[t]{0.76\linewidth}{\raggedright arXiv v2 (2026-07-09)}\\
\textbf{Source record:} \texttt{audit/paper\_statement\_map.json}\\
\textbf{Rows:} 18\\
\textbf{Generated:} 2026-09-06\\
\textbf{Source URL:} \href{https://arxiv.org/abs/2606.08872v2}{official source}





\section*{How to use this packet}
For each source claim, compare the exact source excerpts with the expanded
PaperInterface specification. Mark up this PDF or fill the blank reviewer box in the TeX source;
return the annotated file for reconciliation into the paper-local human-review
log.

\section*{Open the interactive dashboard}
The interactive dashboard is an optional alternative to using this PDF; it
presents the same review surface rather than an additional review requirement.
After forking and cloning (or pulling) AppliedModelingLib, run the following from the
repository root. The cache command is needed only the first time in a new clone.
\begin{ReviewVerbatim}
cd <your-repository-checkout>
lake exe cache get
python3 scripts/review_dashboard.py --paper HT26EFXChores --serve
\end{ReviewVerbatim}
Open \url{http://127.0.0.1:8765/} in a browser and keep that terminal running
while reviewing. The dashboard records saved annotations in this paper's local
reviewer-owned review record.

\section*{Contents}
\begin{itemize}[leftmargin=1.2em]
\item \hyperlink{library-section-de5e5601682a1524}{Semantic library prerequisites (3; not paper claims)}
\begin{itemize}[leftmargin=1.2em]
\item \hyperlink{library-prerequisite-cc4bb81e3f6310d6}{Applied\allowbreak{}Modeling\allowbreak{}Lib.\allowbreak{}Fair\allowbreak{}Division.\allowbreak{}EFX\allowbreak{}For\allowbreak{}Chores}
\item \hyperlink{library-prerequisite-283bac4b270e431c}{Applied\allowbreak{}Modeling\allowbreak{}Lib.\allowbreak{}Fair\allowbreak{}Division.\allowbreak{}Envy\allowbreak{}Free\allowbreak{}For\allowbreak{}Chores}
\item \hyperlink{library-prerequisite-95d7f6c33ea68817}{Applied\allowbreak{}Modeling\allowbreak{}Lib.\allowbreak{}Fair\allowbreak{}Division.\allowbreak{}Pareto\allowbreak{}Optimal\allowbreak{}For\allowbreak{}Chores}
\end{itemize}
\item \hyperlink{paper-prerequisite-section-5ef5ef0364b6939c}{Paper-specific semantic prerequisites (3; not paper claims)}
\begin{itemize}[leftmargin=1.2em]
\item \hyperlink{paper-prerequisite-fd04b7b4fbc6f835}{HT26\allowbreak{}EFX\allowbreak{}Chores.\allowbreak{}canonical\allowbreak{}Short\allowbreak{}Long\allowbreak{}Labels\allowbreak{}Source\allowbreak{}Definition}
\item \hyperlink{paper-prerequisite-efb704bbc61ee065}{HT26\allowbreak{}EFX\allowbreak{}Chores.\allowbreak{}canonical\allowbreak{}Small\allowbreak{}Chore\allowbreak{}Allocation\allowbreak{}Source\allowbreak{}Definition}
\item \hyperlink{paper-prerequisite-b0f133a4a33c9a55}{HT26\allowbreak{}EFX\allowbreak{}Chores.\allowbreak{}super\allowbreak{}Canonical\allowbreak{}Small\allowbreak{}Chore\allowbreak{}Allocation\allowbreak{}Source\allowbreak{}Definition}
\end{itemize}
\item \hyperlink{source-section-0c7af8265958a9ea}{Main-text source claims (16)}
\begin{itemize}[leftmargin=1.2em]
\item \hyperlink{source-claim-707062c38758c055}{Theorem 1 (tri-valued EFX nonexistence)}
\item \hyperlink{source-claim-2321d8010acce683}{Four-agent construction proposition: at most one A item per bundle}
\item \hyperlink{source-claim-120f295a37ffc0ce}{Four-agent construction proposition: no-A bundle is expensive}
\item \hyperlink{source-claim-8bed2d2e9152e71d}{Theorem 2 (EFX and Pareto-optimality incompatibility)}
\item \hyperlink{source-claim-6906c411c024867a}{Theorem 2 explicit EFX construction}
\item \hyperlink{source-claim-5145ed0daed8afc6}{Theorem 2 proposition: every EFX agent receives a large item}
\item \hyperlink{source-claim-8a3c68dd9df286a5}{Theorem 2 proposition: EFX cost lower bounds}
\item \hyperlink{source-claim-03ed4c5388a9d78f}{Theorem 3 (four-agent bi-valued EFX existence)}
\item \hyperlink{source-claim-e126eba4635c1cd5}{Lemma (M34 insertion)}
\item \hyperlink{source-claim-5fbcb5b015031605}{Remark (general-agent insertion)}
\item \hyperlink{source-claim-58fb9c46acb21eca}{Lemma (composition)}
\item \hyperlink{source-claim-1b77c29da35b62c2}{Lemma (canonical allocation properties)}
\item \hyperlink{source-claim-0869bc0832b92b55}{Lemma (balanced orientation)}
\item \hyperlink{source-claim-6263bf063357fb29}{Lemma (EFX allocation of M2)}
\item \hyperlink{source-claim-687ed4b91a799ce3}{Lemma (properties of the M2 EFX allocations)}
\item \hyperlink{source-claim-c475426ec81f3892}{Lemma (exceptional-residue combination)}
\end{itemize}
\item \hyperlink{source-section-1743f0a28d68d6cc}{Appendix and supplement source claims (2)}
\begin{itemize}[leftmargin=1.2em]
\item \hyperlink{source-claim-bf31a0d3fb856868}{Appendix A proposition: no bundle contains two A items}
\item \hyperlink{source-claim-18b6e05d84d21367}{Appendix A proposition: p1 and p2 lower bounds}
\end{itemize}
\end{itemize}

\clearpage
\hypertarget{library-section-de5e5601682a1524}{}
\section*{Material library prerequisites}
\hypertarget{library-prerequisite-cc4bb81e3f6310d6}{}
\subsection*{\small\ttfamily\raggedright Applied\allowbreak{}Modeling\allowbreak{}Lib.\allowbreak{}Fair\allowbreak{}Division.\allowbreak{}EFX\allowbreak{}For\allowbreak{}Chores}
\textbf{Source locator:} {\footnotesize\raggedright cited publication, lines 260-\allowbreak{}275\par}
\paragraph{Verbatim paper-source connection}
\begin{ReviewVerbatim}
\begin{definition}[EFX for chores]
An allocation $X=(X_1,\ldots,X_n)$ is \emph{EFX} if for every pair of agents $i,j$ we have either
\begin{itemize}
    \item $X_i=\emptyset$, or
    \item for every item $g\in X_i$, $c_i(X_i\setminus\{g\})\le c_i(X_j).$
\end{itemize}
Equivalently, if
\[
    \tau_i(X):=\tau_i(X_i):=
    \begin{cases}
        0, & X_i=\emptyset,\\
        c_i(X_i)-\min_{g\in X_i}c_i(g), & X_i\ne\emptyset,
    \end{cases}
\]
then $X$ is EFX if and only if $\tau_i(X)\le c_i(X_j)$ for all $i,j$.
\end{definition}

[Next verbatim source excerpt]

We will use ``item'' and ``chore'' interchangeably in this paper to refer to an element in $M$.
Each agent $i$'s cost function $c_i:2^M\to\mathbb{R}_{\geq0}$ is additive:
$$c_i(S)=\sum_{g\in S}c_i(\{g\}).$$
We write $c_i(g_j)$ in short for $c_i(\{g_j\})$ for notation simplicity.
We will use ``cost function'' and ``disutility function'' interchangeably in this paper.
\end{ReviewVerbatim}



\paragraph{Lean-expanded library semantic target}
\begin{ReviewVerbatim}
type:
{Agent : Type u_1} →
  {Item : Type u_2} →
    [DecidableEq Item] →
      (Agent → AppliedModelingLib.FairDivision.Bundle Item → ℝ) →
        AppliedModelingLib.FairDivision.Allocation Agent Item → Prop

value:
fun {Agent : Type u_1} {Item : Type u_2} [DecidableEq Item]
    (cost : Agent → AppliedModelingLib.FairDivision.Bundle Item → ℝ)
    (allocation : AppliedModelingLib.FairDivision.Allocation Agent Item) =>
  ∀ (i j : Agent), allocation i = ∅ ∨ ∀ item ∈ allocation i, cost i (allocation i \ {item}) ≤ cost i (allocation j)
\end{ReviewVerbatim}

\paragraph{Exact Lean library declaration}
\begin{ReviewVerbatim}
/-- Envy-freeness up to any item for chores: after removing any item from an
agent's own bundle, that agent incurs no more cost than for another bundle. -/
def EFXForChores [DecidableEq Item] (cost : Agent → Bundle Item → ℝ)
    (allocation : Allocation Agent Item) : Prop :=
  ∀ i j, allocation i = ∅ ∨
    ∀ item ∈ allocation i, cost i (allocation i \ {item}) ≤ cost i (allocation j)
\end{ReviewVerbatim}

\paragraph{Recorded source-to-library screening}
\textbf{Verdict:} \texttt{matches}
\\
{\raggedright
\textbf{Reason:} The target states that every agent pair either has an empty own bundle or satisfies the source all-\allowbreak{}item removal inequality.\allowbreak{}
\par}
\reviewmatch{Matches source input}
\noindent\textbf{Reviewer annotation}\par
\reviewerbox
\clearpage
\hypertarget{library-prerequisite-283bac4b270e431c}{}
\subsection*{\small\ttfamily\raggedright Applied\allowbreak{}Modeling\allowbreak{}Lib.\allowbreak{}Fair\allowbreak{}Division.\allowbreak{}Envy\allowbreak{}Free\allowbreak{}For\allowbreak{}Chores}
\textbf{Source locator:} {\footnotesize\raggedright cited publication, lines 253-\allowbreak{}255\par}
\paragraph{Verbatim paper-source connection}
\begin{ReviewVerbatim}
\begin{definition}[EF for chores]
An allocation $X=(X_1,\ldots,X_n)$ is \emph{envy-free} if for every pair of agents $i,j$ we have $c_i(X_i)\leq c_i(X_j)$.
\end{definition}
\end{ReviewVerbatim}



\paragraph{Lean-expanded library semantic target}
\begin{ReviewVerbatim}
type:
{Agent : Type u_1} →
  {Item : Type u_2} →
    (Agent → AppliedModelingLib.FairDivision.Bundle Item → ℝ) →
      AppliedModelingLib.FairDivision.Allocation Agent Item → Prop

value:
fun {Agent : Type u_1} {Item : Type u_2} (cost : Agent → AppliedModelingLib.FairDivision.Bundle Item → ℝ)
    (allocation : AppliedModelingLib.FairDivision.Allocation Agent Item) =>
  ∀ (i j : Agent), cost i (allocation i) ≤ cost i (allocation j)
\end{ReviewVerbatim}

\paragraph{Exact Lean library declaration}
\begin{ReviewVerbatim}
/-- An allocation is envy-free for chores when each agent's own cost is no
greater than their cost for any other agent's bundle. -/
def EnvyFreeForChores (cost : Agent → Bundle Item → ℝ)
    (allocation : Allocation Agent Item) : Prop :=
  ∀ i j, cost i (allocation i) ≤ cost i (allocation j)
\end{ReviewVerbatim}

\paragraph{Recorded source-to-library screening}
\textbf{Verdict:} \texttt{matches}
\\
{\raggedright
\textbf{Reason:} The target states exactly that every agent's own bundle cost is at most that agent's cost for every other bundle.\allowbreak{}
\par}
\reviewmatch{Matches source input}
\noindent\textbf{Reviewer annotation}\par
\reviewerbox
\clearpage
\hypertarget{library-prerequisite-95d7f6c33ea68817}{}
\subsection*{\small\ttfamily\raggedright Applied\allowbreak{}Modeling\allowbreak{}Lib.\allowbreak{}Fair\allowbreak{}Division.\allowbreak{}Pareto\allowbreak{}Optimal\allowbreak{}For\allowbreak{}Chores}
\textbf{Source locator:} {\footnotesize\raggedright cited publication, lines 280-\allowbreak{}287\par}
\paragraph{Verbatim paper-source connection}
\begin{ReviewVerbatim}
\begin{definition}[Pareto-Optimality]
    An allocation $Y=(Y_1,\ldots,Y_n)$ \emph{Pareto-dominates} an allocation $X=(X_1,\ldots,X_n)$ if
    \begin{itemize}
        \item $c_i(Y_i)\leq c_i(X_i)$ holds for all $i\in\{1,\ldots,n\}$, and
        \item there exists $i\in\{1,\ldots,n\},\;$ $c_i(Y_i)<c_i(X_i)$.
    \end{itemize}
    An allocation $X=(X_1,\ldots,X_n)$ is \emph{Pareto-optimal} if it is not Pareto-dominated by any allocation.
\end{definition}
\end{ReviewVerbatim}



\paragraph{Lean-expanded library semantic target}
\begin{ReviewVerbatim}
type:
{Agent : Type u_1} →
  {Item : Type u_2} →
    [DecidableEq Item] →
      (Agent → AppliedModelingLib.FairDivision.Bundle Item → ℝ) →
        Finset Item → AppliedModelingLib.FairDivision.Allocation Agent Item → Prop

value:
fun {Agent : Type u_1} {Item : Type u_2} [DecidableEq Item]
    (cost : Agent → AppliedModelingLib.FairDivision.Bundle Item → ℝ) (chores : Finset Item)
    (allocation : AppliedModelingLib.FairDivision.Allocation Agent Item) =>
  AppliedModelingLib.FairDivision.IsAllocationOf allocation chores ∧
    ¬∃ (improvement : AppliedModelingLib.FairDivision.Allocation Agent Item),
        AppliedModelingLib.FairDivision.IsAllocationOf improvement chores ∧
          AppliedModelingLib.FairDivision.ParetoDominatesForChores cost allocation improvement
\end{ReviewVerbatim}

\paragraph{Exact Lean library declaration}
\begin{ReviewVerbatim}
/-- A complete chore allocation is Pareto-optimal if no other complete
allocation Pareto-dominates it in the cost direction. -/
def ParetoOptimalForChores [DecidableEq Item] (cost : Agent → Bundle Item → ℝ)
    (chores : Finset Item) (allocation : Allocation Agent Item) : Prop :=
  IsAllocationOf allocation chores ∧
    ¬ ∃ improvement : Allocation Agent Item,
      IsAllocationOf improvement chores ∧
        ParetoDominatesForChores cost allocation improvement
\end{ReviewVerbatim}

\paragraph{Recorded source-to-library screening}
\textbf{Verdict:} \texttt{matches}
\\
{\raggedright
\textbf{Reason:} The surrounding preliminaries fix a chore set allocated to the agents, so the source's candidate and comparator allocations are complete allocations of that same chore set, as encoded by IsAllocationOf.\allowbreak{}
\par}
\reviewmatch{Matches source input}
\noindent\textbf{Reviewer annotation}\par
\reviewerbox
\clearpage
\hypertarget{paper-prerequisite-section-5ef5ef0364b6939c}{}
\section*{Paper-specific semantic prerequisites}
\hypertarget{paper-prerequisite-fd04b7b4fbc6f835}{}
\subsection*{\small\ttfamily\raggedright HT26\allowbreak{}EFX\allowbreak{}Chores.\allowbreak{}canonical\allowbreak{}Short\allowbreak{}Long\allowbreak{}Labels\allowbreak{}Source\allowbreak{}Definition}
\noindent\textbf{Paper-local declaration:}\par
{\footnotesize\ttfamily\raggedright HT26\allowbreak{}EFX\allowbreak{}Chores.\allowbreak{}canonical\allowbreak{}Short\allowbreak{}Long\allowbreak{}Labels\allowbreak{}Source\allowbreak{}Definition\par}
\textbf{Source locator:} {\footnotesize\raggedright cited publication, lines 753-\allowbreak{}754\par}
\paragraph{Verbatim paper-source connection}
\begin{ReviewVerbatim}
In a canonical allocation with $b>0$, agents receiving $a$ items are called \emph{short} agents, and agents receiving $a+1$ items are called \emph{long} agents.
    For $b=0$ and $a>0$, every agent in a canonical is both long and short.

[Next verbatim source excerpt]

We will assume $r>2$ in the remaining part of this section.
For a chore $g$, let
\[
 S(g)=\{i\in N:c_i(g)=1\}
\]
be the set of agents for whom $g$ is small.  We partition the chores into
\[
M_{01}=\{g:|S(g)|\le 1\},\qquad
M_2=\{g:|S(g)|=2\},\qquad
M_{34}=\{g:|S(g)|\ge 3\}.
\]
For an item in $M_2$, we write $(i,j)$ for an item small exactly for agents
$i$ and $j$.  Thus the items in $M_2$ form a multigraph on the vertex set
$N$, where an edge $(i,j)$ is a chore small for precisely agents $i$ and $j$.

[Next verbatim source excerpt]

For each agent $i$, let
\[
 n_i=\left|\{g\in M_{01}:S(g)=\{i\}\}\right|
\]
be the number of $M_{01}$ items uniquely small for $i$.  Items in $M_{01}$
with $S(g)=\emptyset$ are all-large items.
Let $|M_{01}|=4a+b$ with $b\in\{0,1,2,3\}$.
As mentioned before, we will make sure each agent receives either $a$ items or $a+1$ items.

Fix quotas $q_i\in\{a,a+1\}$ with $\sum_i q_i=|M_{01}|$.
We consider allocations of $M_{01}$ with quotas $q=(q_1,q_2,q_3,q_4)$ in which each agent $i$ first receives as many of her uniquely-small items as possible, up to quota $q_i$, and the remaining slots are filled with the remaining items.
We will show that all such allocations satisfy EFX.
It is also important that we have the freedom to choose which agents receive $a$ items and which agents receive $a+1$ items.
This freedom enables us to append the allocation of $M_2$ after the allocation of $M_{01}$.

We begin by defining this type of allocations.

\begin{definition}
    An allocation $P=(P_1,P_2,P_3,P_4)$ of $M_{01}$ is \emph{canonical} if, for each agent $i$, $|P_i|\in\{a,a+1\}$ and $P_i$ contains $\min\{|P_i|, n_i\}$ items that are small for agent $i$.

[Next verbatim source excerpt]

be the set of agents for whom $g$ is small.  We partition the chores into
\[
M_{01}=\{g:|S(g)|\le 1\},\qquad
M_2=\{g:|S(g)|=2\},\qquad
M_{34}=\{g:|S(g)|\ge 3\}.
\]
For an item in $M_2$, we write $(i,j)$ for an item small exactly for agents
$i$ and $j$.  Thus the items in $M_2$ form a multigraph on the vertex set
$N$, where an edge $(i,j)$ is a chore small for precisely agents $i$ and $j$.
\end{ReviewVerbatim}



\paragraph{Lean-elaborated paper signature}
\begin{ReviewVerbatim}
type:
(Item : Type) →
  AppliedModelingLib.FairDivision.ChoreCost (Fin 4) Item →
    Finset Item →
      ℕ →
        ℕ →
          (Fin 4 → ℕ) → AppliedModelingLib.FairDivision.Allocation (Fin 4) Item → Finset (Fin 4) → Finset (Fin 4) → Prop

value:
fun (Item : Type) (cost : AppliedModelingLib.FairDivision.ChoreCost (Fin 4) Item) (chores : Finset Item) (a b : ℕ)
    (quota : Fin 4 → ℕ) (allocation : AppliedModelingLib.FairDivision.Allocation (Fin 4) Item)
    (shortAgents longAgents : Finset (Fin 4)) =>
  HT26EFXChores.canonicalSmallChoreAllocationSourceDefinition Item cost chores a b quota allocation ∧
    (0 < b →
        (∀ (agent : Fin 4), agent ∈ shortAgents ↔ quota agent = a) ∧
          ∀ (agent : Fin 4), agent ∈ longAgents ↔ quota agent = a + 1) ∧
      (b = 0 → 0 < a → shortAgents = Finset.univ ∧ longAgents = Finset.univ)
\end{ReviewVerbatim}


\paragraph{Recorded source-to-declaration screening}
\textbf{Verdict:} \texttt{matches}
\\
{\raggedright
\textbf{Reason:} The target identifies short and long agents by the source quota convention, including the b=0, a>0 convention and the source-\allowbreak{}undefined b=0, a=0 case.\allowbreak{}
\par}
\reviewmatch{Matches source input}
\noindent\textbf{Reviewer annotation}\par
\reviewerbox
\clearpage
\hypertarget{paper-prerequisite-efb704bbc61ee065}{}
\subsection*{\small\ttfamily\raggedright HT26\allowbreak{}EFX\allowbreak{}Chores.\allowbreak{}canonical\allowbreak{}Small\allowbreak{}Chore\allowbreak{}Allocation\allowbreak{}Source\allowbreak{}Definition}
\noindent\textbf{Paper-local declaration:}\par
{\footnotesize\ttfamily\raggedright HT26\allowbreak{}EFX\allowbreak{}Chores.\allowbreak{}canonical\allowbreak{}Small\allowbreak{}Chore\allowbreak{}Allocation\allowbreak{}Source\allowbreak{}Definition\par}
\textbf{Source locator:} {\footnotesize\raggedright cited publication, lines 752\par}
\paragraph{Verbatim paper-source connection}
\begin{ReviewVerbatim}
An allocation $P=(P_1,P_2,P_3,P_4)$ of $M_{01}$ is \emph{canonical} if, for each agent $i$, $|P_i|\in\{a,a+1\}$ and $P_i$ contains $\min\{|P_i|, n_i\}$ items that are small for agent $i$.

[Next verbatim source excerpt]

We will assume $r>2$ in the remaining part of this section.
For a chore $g$, let
\[
 S(g)=\{i\in N:c_i(g)=1\}
\]
be the set of agents for whom $g$ is small.  We partition the chores into
\[
M_{01}=\{g:|S(g)|\le 1\},\qquad
M_2=\{g:|S(g)|=2\},\qquad
M_{34}=\{g:|S(g)|\ge 3\}.
\]
For an item in $M_2$, we write $(i,j)$ for an item small exactly for agents
$i$ and $j$.  Thus the items in $M_2$ form a multigraph on the vertex set
$N$, where an edge $(i,j)$ is a chore small for precisely agents $i$ and $j$.

[Next verbatim source excerpt]

For each agent $i$, let
\[
 n_i=\left|\{g\in M_{01}:S(g)=\{i\}\}\right|
\]
be the number of $M_{01}$ items uniquely small for $i$.  Items in $M_{01}$
with $S(g)=\emptyset$ are all-large items.
Let $|M_{01}|=4a+b$ with $b\in\{0,1,2,3\}$.
As mentioned before, we will make sure each agent receives either $a$ items or $a+1$ items.

Fix quotas $q_i\in\{a,a+1\}$ with $\sum_i q_i=|M_{01}|$.
We consider allocations of $M_{01}$ with quotas $q=(q_1,q_2,q_3,q_4)$ in which each agent $i$ first receives as many of her uniquely-small items as possible, up to quota $q_i$, and the remaining slots are filled with the remaining items.
We will show that all such allocations satisfy EFX.
It is also important that we have the freedom to choose which agents receive $a$ items and which agents receive $a+1$ items.
This freedom enables us to append the allocation of $M_2$ after the allocation of $M_{01}$.

We begin by defining this type of allocations.


[Next verbatim source excerpt]

be the set of agents for whom $g$ is small.  We partition the chores into
\[
M_{01}=\{g:|S(g)|\le 1\},\qquad
M_2=\{g:|S(g)|=2\},\qquad
M_{34}=\{g:|S(g)|\ge 3\}.
\]
For an item in $M_2$, we write $(i,j)$ for an item small exactly for agents
$i$ and $j$.  Thus the items in $M_2$ form a multigraph on the vertex set
$N$, where an edge $(i,j)$ is a chore small for precisely agents $i$ and $j$.
\end{ReviewVerbatim}



\paragraph{Lean-elaborated paper signature}
\begin{ReviewVerbatim}
type:
(Item : Type) →
  AppliedModelingLib.FairDivision.ChoreCost (Fin 4) Item →
    Finset Item → ℕ → ℕ → (Fin 4 → ℕ) → AppliedModelingLib.FairDivision.Allocation (Fin 4) Item → Prop

value:
fun (Item : Type) (cost : AppliedModelingLib.FairDivision.ChoreCost (Fin 4) Item) (chores : Finset Item) (a b : ℕ)
    (quota : Fin 4 → ℕ) (allocation : AppliedModelingLib.FairDivision.Allocation (Fin 4) Item) =>
  b ≤ 3 ∧
    chores.card = 4 * a + b ∧
      (∀ item ∈ chores, AppliedModelingLib.FairDivision.IsSmallForAtMostOne cost item) ∧
        (∀ (agent : Fin 4), quota agent = a ∨ quota agent = a + 1) ∧
          Finset.univ.sum quota = chores.card ∧
            AppliedModelingLib.FairDivision.IsAllocationOf allocation chores ∧
              ∀ (agent : Fin 4),
                Finset.card (allocation agent) = quota agent ∧
                  (AppliedModelingLib.FairDivision.ownSmallChoreSet cost (allocation agent) agent).card =
                    min (quota agent) (AppliedModelingLib.FairDivision.ownSmallChoreSet cost chores agent).card
\end{ReviewVerbatim}


\paragraph{Recorded source-to-declaration screening}
\textbf{Verdict:} \texttt{matches}
\\
{\raggedright
\textbf{Reason:} The target encodes b at most three, the 4a+b chore count, balanced quotas, complete allocation, and exactly the source maximum number of own-\allowbreak{}small chores per bundle.\allowbreak{}
\par}
\reviewmatch{Matches source input}
\noindent\textbf{Reviewer annotation}\par
\reviewerbox
\clearpage
\hypertarget{paper-prerequisite-b0f133a4a33c9a55}{}
\subsection*{\small\ttfamily\raggedright HT26\allowbreak{}EFX\allowbreak{}Chores.\allowbreak{}super\allowbreak{}Canonical\allowbreak{}Small\allowbreak{}Chore\allowbreak{}Allocation\allowbreak{}Source\allowbreak{}Definition}
\noindent\textbf{Paper-local declaration:}\par
{\footnotesize\ttfamily\raggedright HT26\allowbreak{}EFX\allowbreak{}Chores.\allowbreak{}super\allowbreak{}Canonical\allowbreak{}Small\allowbreak{}Chore\allowbreak{}Allocation\allowbreak{}Source\allowbreak{}Definition\par}
\textbf{Source locator:} {\footnotesize\raggedright cited publication, lines 755\par}
\paragraph{Verbatim paper-source connection}
\begin{ReviewVerbatim}
For $b>0$, an allocation $P=(P_1,P_2,P_3,P_4)$ is \emph{super-canonical} if it is canonical and satisfies $c_i(P_i)\leq c_i(P_j)-r$ for every short agent $i$ and every long agent $j$.

[Next verbatim source excerpt]

We will assume $r>2$ in the remaining part of this section.
For a chore $g$, let
\[
 S(g)=\{i\in N:c_i(g)=1\}
\]
be the set of agents for whom $g$ is small.  We partition the chores into
\[
M_{01}=\{g:|S(g)|\le 1\},\qquad
M_2=\{g:|S(g)|=2\},\qquad
M_{34}=\{g:|S(g)|\ge 3\}.
\]
For an item in $M_2$, we write $(i,j)$ for an item small exactly for agents
$i$ and $j$.  Thus the items in $M_2$ form a multigraph on the vertex set
$N$, where an edge $(i,j)$ is a chore small for precisely agents $i$ and $j$.

[Next verbatim source excerpt]

For each agent $i$, let
\[
 n_i=\left|\{g\in M_{01}:S(g)=\{i\}\}\right|
\]
be the number of $M_{01}$ items uniquely small for $i$.  Items in $M_{01}$
with $S(g)=\emptyset$ are all-large items.
Let $|M_{01}|=4a+b$ with $b\in\{0,1,2,3\}$.
As mentioned before, we will make sure each agent receives either $a$ items or $a+1$ items.

Fix quotas $q_i\in\{a,a+1\}$ with $\sum_i q_i=|M_{01}|$.
We consider allocations of $M_{01}$ with quotas $q=(q_1,q_2,q_3,q_4)$ in which each agent $i$ first receives as many of her uniquely-small items as possible, up to quota $q_i$, and the remaining slots are filled with the remaining items.
We will show that all such allocations satisfy EFX.
It is also important that we have the freedom to choose which agents receive $a$ items and which agents receive $a+1$ items.
This freedom enables us to append the allocation of $M_2$ after the allocation of $M_{01}$.

We begin by defining this type of allocations.


[Next verbatim source excerpt]

\begin{definition}
    An allocation $P=(P_1,P_2,P_3,P_4)$ of $M_{01}$ is \emph{canonical} if, for each agent $i$, $|P_i|\in\{a,a+1\}$ and $P_i$ contains $\min\{|P_i|, n_i\}$ items that are small for agent $i$.
    In a canonical allocation with $b>0$, agents receiving $a$ items are called \emph{short} agents, and agents receiving $a+1$ items are called \emph{long} agents.
    For $b=0$ and $a>0$, every agent in a canonical is both long and short.

[Next verbatim source excerpt]

be the set of agents for whom $g$ is small.  We partition the chores into
\[
M_{01}=\{g:|S(g)|\le 1\},\qquad
M_2=\{g:|S(g)|=2\},\qquad
M_{34}=\{g:|S(g)|\ge 3\}.
\]
For an item in $M_2$, we write $(i,j)$ for an item small exactly for agents
$i$ and $j$.  Thus the items in $M_2$ form a multigraph on the vertex set
$N$, where an edge $(i,j)$ is a chore small for precisely agents $i$ and $j$.
\end{ReviewVerbatim}



\paragraph{Lean-elaborated paper signature}
\begin{ReviewVerbatim}
type:
(Item : Type) →
  ℝ →
    AppliedModelingLib.FairDivision.ChoreCost (Fin 4) Item →
      Finset Item → ℕ → ℕ → (Fin 4 → ℕ) → AppliedModelingLib.FairDivision.Allocation (Fin 4) Item → Prop

value:
fun (Item : Type) (r : ℝ) (cost : AppliedModelingLib.FairDivision.ChoreCost (Fin 4) Item) (chores : Finset Item)
    (a b : ℕ) (quota : Fin 4 → ℕ) (allocation : AppliedModelingLib.FairDivision.Allocation (Fin 4) Item) =>
  0 < b ∧
    ∃ (shortAgents : Finset (Fin 4)) (longAgents : Finset (Fin 4)),
      HT26EFXChores.canonicalShortLongLabelsSourceDefinition Item cost chores a b quota allocation shortAgents
          longAgents ∧
        ∀ shortAgent ∈ shortAgents,
          ∀ longAgent ∈ longAgents,
            AppliedModelingLib.FairDivision.additiveChoreCost cost shortAgent (allocation shortAgent) ≤
              AppliedModelingLib.FairDivision.additiveChoreCost cost shortAgent (allocation longAgent) - r
\end{ReviewVerbatim}


\paragraph{Recorded source-to-declaration screening}
\textbf{Verdict:} \texttt{matches}
\\
{\raggedright
\textbf{Reason:} The target requires b positive, canonicality with the source short/\allowbreak{}long convention, and exactly the stated r cost gap from every short agent to every long agent.\allowbreak{}
\par}
\reviewmatch{Matches source input}
\noindent\textbf{Reviewer annotation}\par
\reviewerbox

\clearpage
\hypertarget{source-claim-707062c38758c055}{}
\hypertarget{source-section-0c7af8265958a9ea}{}
\section*{Main-text source claims}
\subsection*{Review row 1}
\textbf{Source locator:} {\footnotesize\raggedright cited publication, lines 304-\allowbreak{}306\par}
\paragraph{Verbatim source input}
\begin{ReviewVerbatim}
\begin{theorem}\label{thm:tri}
    For any $n\geq 4$, there exists a tri-valued instance with $n$ agents where no EFX allocation exists.
\end{theorem}

[Next verbatim source excerpt]

\section{Preliminaries}
A set $M$ of $m$ chores $g_1,\ldots,g_m$ is allocated to a set $N$ of $n$ agents $1,\ldots,n$.
We will use ``item'' and ``chore'' interchangeably in this paper to refer to an element in $M$.
Each agent $i$'s cost function $c_i:2^M\to\mathbb{R}_{\geq0}$ is additive:
$$c_i(S)=\sum_{g\in S}c_i(\{g\}).$$
We write $c_i(g_j)$ in short for $c_i(\{g_j\})$ for notation simplicity.
We will use ``cost function'' and ``disutility function'' interchangeably in this paper.

[Next verbatim source excerpt]

We say that the cost function is tri-valued if there exist constants $p,q,r\geq 0$ such that $c_i(g_j)\in\{p,q,r\}$ for any $i\in N$ and $g_j\in M$.
The cost function is bi-valued if there exist constants $p,q\geq 0$ such that $c_i(g_j)\in\{p,q\}$ for any $i\in N$ and $g_j\in M$.
For bi-valued cost functions, we will focus on the case $0<p<q$ (the case $0=p<q$ corresponds to the setting with binary cost functions and we know an allocation that is EFX and Pareto-optimal always exists), and we will assume $c_i(g_j)\in\{1,r\}$ (for any $i\in N$ and $g_j\in M$) instead.
We say that the item $g_j$ is ``large'' for agent $i$ if $c_i(g_j)=r$ and it is ``small'' for agent $i$ if $c_i(g_j)=1$.
\end{ReviewVerbatim}



\paragraph{Expanded PaperInterface specification}
\begin{ReviewVerbatim}
∀ (n : ℕ),
  4 ≤ n →
    ∃ (m : ℕ) (choreInstance : AppliedModelingLib.FairDivision.AdditiveChoreInstance (Fin n) (Fin m)),
      AppliedModelingLib.FairDivision.IsTriValuedChoreCost choreInstance.cost ∧
        ∀ (allocation : AppliedModelingLib.FairDivision.Allocation (Fin n) (Fin m)),
          choreInstance.IsFeasible allocation → ¬choreInstance.IsEFX allocation
\end{ReviewVerbatim}



\paragraph{Lean proof endpoint (built separately)}\begin{ReviewVerbatim}
HT26EFXChores.tri_valued_nonexistence
\end{ReviewVerbatim}

\paragraph{Recorded source-to-Spec screening}
\textbf{Verdict:} \texttt{matches}
\\
{\raggedright
\textbf{Reason:} Exact n-\allowbreak{}at-\allowbreak{}least-\allowbreak{}four existential tri-\allowbreak{}valued additive nonnegative instance with no feasible EFX allocation.\allowbreak{}
\par}
\reviewmatch{Matches source input}
\noindent\textbf{Reviewer annotation}\par
\reviewerbox
\clearpage
\hypertarget{source-claim-2321d8010acce683}{}
\section*{Review row 2}
\textbf{Source locator:} {\footnotesize\raggedright cited publication, lines 340-\allowbreak{}342\par}
\paragraph{Verbatim source input}
\begin{ReviewVerbatim}
\begin{proposition}
    No agent can receive more than one item in $A$.
\end{proposition}

[Next verbatim source excerpt]

The instance consists of $4$ agents and $13$ items.
The items are partitioned into three groups:
$$A=\{a_1,a_2,a_3\},\quad B=\{b_1,b_2,b_3,b_4,b_5\},\quad\mbox{and}\quad C=\{c_1,c_2,c_3,c_4,c_5\}.$$
The items group $A$ consists of ``large'' items where each agent values each of them at $20$.
For each item in $B$, agents $1$ and $2$ value it at $1$ and agents $3$ and $4$ value it at $7$.
For each item in $C$, agents $1$ and $2$ value it at $7$ and agents $3$ and $4$ value it at $1$.

[Next verbatim source excerpt]

In the remaining part of this section, we will show that no EFX allocation exists in this instance.
Suppose for contradiction an EFX allocation $X=(X_1,X_2,X_3,X_4)$ exists.
Firstly, we will show that no one can receive more than one large item in $A$.

[Next verbatim source excerpt]

\begin{definition}[EFX for chores]
An allocation $X=(X_1,\ldots,X_n)$ is \emph{EFX} if for every pair of agents $i,j$ we have either
\begin{itemize}
    \item $X_i=\emptyset$, or
    \item for every item $g\in X_i$, $c_i(X_i\setminus\{g\})\le c_i(X_j).$
\end{itemize}
Equivalently, if
\[
    \tau_i(X):=\tau_i(X_i):=
    \begin{cases}
        0, & X_i=\emptyset,\\
        c_i(X_i)-\min_{g\in X_i}c_i(g), & X_i\ne\emptyset,
    \end{cases}
\]
then $X$ is EFX if and only if $\tau_i(X)\le c_i(X_j)$ for all $i,j$.
\end{definition}
\end{ReviewVerbatim}



\paragraph{Expanded PaperInterface specification}
\begin{ReviewVerbatim}
∀ (allocation : AppliedModelingLib.FairDivision.Allocation (Fin 4) (Fin 13)),
  AppliedModelingLib.FairDivision.IsAllocationOf allocation Finset.univ →
    AppliedModelingLib.FairDivision.EFXForChores
        (AppliedModelingLib.FairDivision.additiveChoreCost fun (agent : Fin 4) (item : Fin 13) =>
          if ↑item < 3 then 20 else if ↑agent < 2 then if ↑item < 8 then 1 else 7 else if ↑item < 8 then 7 else 1)
        allocation →
      ∀ (agent : Fin 4), {item ∈ allocation agent | ↑item < 3}.card ≤ 1
\end{ReviewVerbatim}



\paragraph{Lean proof endpoint (built separately)}\begin{ReviewVerbatim}
HT26EFXChores.tri_four_no_two_large
\end{ReviewVerbatim}

\paragraph{Recorded source-to-Spec screening}
\textbf{Verdict:} \texttt{matches}
\\
{\raggedright
\textbf{Reason:} Exact concrete four-\allowbreak{}agent, 13-\allowbreak{}item instance and at-\allowbreak{}most-\allowbreak{}one-\allowbreak{}A-\allowbreak{}item EFX consequence.\allowbreak{}
\par}
\reviewmatch{Matches source input}
\noindent\textbf{Reviewer annotation}\par
\reviewerbox
\clearpage
\hypertarget{source-claim-120f295a37ffc0ce}{}
\section*{Review row 3}
\textbf{Source locator:} {\footnotesize\raggedright cited publication, lines 363-\allowbreak{}365\par}
\paragraph{Verbatim source input}
\begin{ReviewVerbatim}
\begin{proposition}
    For any agent $i$, we have $c_i(X_1)\geq 20$.
\end{proposition}

[Next verbatim source excerpt]

The instance consists of $4$ agents and $13$ items.
The items are partitioned into three groups:
$$A=\{a_1,a_2,a_3\},\quad B=\{b_1,b_2,b_3,b_4,b_5\},\quad\mbox{and}\quad C=\{c_1,c_2,c_3,c_4,c_5\}.$$
The items group $A$ consists of ``large'' items where each agent values each of them at $20$.
For each item in $B$, agents $1$ and $2$ value it at $1$ and agents $3$ and $4$ value it at $7$.
For each item in $C$, agents $1$ and $2$ value it at $7$ and agents $3$ and $4$ value it at $1$.

[Next verbatim source excerpt]

In the remaining part of this section, we will show that no EFX allocation exists in this instance.
Suppose for contradiction an EFX allocation $X=(X_1,X_2,X_3,X_4)$ exists.
Firstly, we will show that no one can receive more than one large item in $A$.

[Next verbatim source excerpt]

With the above proposition, it must be that one agent does not receive any large item in $A$ and each of the remaining three agents receives exactly one item in $A$.
Suppose without loss of generality that agent $1$ does not receive any item from $A$.

[Next verbatim source excerpt]

\begin{definition}[EFX for chores]
An allocation $X=(X_1,\ldots,X_n)$ is \emph{EFX} if for every pair of agents $i,j$ we have either
\begin{itemize}
    \item $X_i=\emptyset$, or
    \item for every item $g\in X_i$, $c_i(X_i\setminus\{g\})\le c_i(X_j).$
\end{itemize}
Equivalently, if
\[
    \tau_i(X):=\tau_i(X_i):=
    \begin{cases}
        0, & X_i=\emptyset,\\
        c_i(X_i)-\min_{g\in X_i}c_i(g), & X_i\ne\emptyset,
    \end{cases}
\]
then $X$ is EFX if and only if $\tau_i(X)\le c_i(X_j)$ for all $i,j$.
\end{definition}
\end{ReviewVerbatim}



\paragraph{Expanded PaperInterface specification}
\begin{ReviewVerbatim}
∀ (allocation : AppliedModelingLib.FairDivision.Allocation (Fin 4) (Fin 13)),
  AppliedModelingLib.FairDivision.IsAllocationOf allocation Finset.univ →
    AppliedModelingLib.FairDivision.EFXForChores
        (AppliedModelingLib.FairDivision.additiveChoreCost fun (agent : Fin 4) (item : Fin 13) =>
          if ↑item < 3 then 20 else if ↑agent < 2 then if ↑item < 8 then 1 else 7 else if ↑item < 8 then 7 else 1)
        allocation →
      {item ∈ allocation 0 | ↑item < 3}.card = 0 →
        (∀ (agent : Fin 4), agent ≠ 0 → {item ∈ allocation agent | ↑item < 3}.card = 1) →
          ∀ (agent : Fin 4),
            20 ≤
              AppliedModelingLib.FairDivision.additiveChoreCost
                (fun (agent : Fin 4) (item : Fin 13) =>
                  if ↑item < 3 then 20
                  else if ↑agent < 2 then if ↑item < 8 then 1 else 7 else if ↑item < 8 then 7 else 1)
                agent (allocation 0)
\end{ReviewVerbatim}



\paragraph{Lean proof endpoint (built separately)}\begin{ReviewVerbatim}
HT26EFXChores.tri_four_no_a_bundle_expensive
\end{ReviewVerbatim}

\paragraph{Recorded source-to-Spec screening}
\textbf{Verdict:} \texttt{matches}
\\
{\raggedright
\textbf{Reason:} Under the supplied WLOG distribution of the three A-\allowbreak{}items, asserts every agent values agent 1's A-\allowbreak{}free bundle at least 20.\allowbreak{}
\par}
\reviewmatch{Matches source input}
\noindent\textbf{Reviewer annotation}\par
\reviewerbox
\clearpage
\hypertarget{source-claim-8bed2d2e9152e71d}{}
\section*{Review row 4}
\textbf{Source locator:} {\footnotesize\raggedright cited publication, lines 416-\allowbreak{}418\par}
\paragraph{Verbatim source input}
\begin{ReviewVerbatim}
\begin{theorem}\label{thm:EFXPO}
    For every $n\geq 4$ and $r>\lceil n/2\rceil+1$, there exists a $(1,r)$-bi-valued instance with $n$ agents where every EFX allocation is not Pareto-optimal.
\end{theorem}

[Next verbatim source excerpt]

\section{Preliminaries}
A set $M$ of $m$ chores $g_1,\ldots,g_m$ is allocated to a set $N$ of $n$ agents $1,\ldots,n$.
We will use ``item'' and ``chore'' interchangeably in this paper to refer to an element in $M$.
Each agent $i$'s cost function $c_i:2^M\to\mathbb{R}_{\geq0}$ is additive:
$$c_i(S)=\sum_{g\in S}c_i(\{g\}).$$
We write $c_i(g_j)$ in short for $c_i(\{g_j\})$ for notation simplicity.
We will use ``cost function'' and ``disutility function'' interchangeably in this paper.

[Next verbatim source excerpt]

We say that the cost function is tri-valued if there exist constants $p,q,r\geq 0$ such that $c_i(g_j)\in\{p,q,r\}$ for any $i\in N$ and $g_j\in M$.
The cost function is bi-valued if there exist constants $p,q\geq 0$ such that $c_i(g_j)\in\{p,q\}$ for any $i\in N$ and $g_j\in M$.
For bi-valued cost functions, we will focus on the case $0<p<q$ (the case $0=p<q$ corresponds to the setting with binary cost functions and we know an allocation that is EFX and Pareto-optimal always exists), and we will assume $c_i(g_j)\in\{1,r\}$ (for any $i\in N$ and $g_j\in M$) instead.
We say that the item $g_j$ is ``large'' for agent $i$ if $c_i(g_j)=r$ and it is ``small'' for agent $i$ if $c_i(g_j)=1$.
\end{ReviewVerbatim}



\paragraph{Expanded PaperInterface specification}
\begin{ReviewVerbatim}
∀ (n : ℕ),
  4 ≤ n →
    ∀ (r : ℝ),
      ↑((n + 1) / 2) + 1 < r →
        ∃ (m : ℕ) (choreInstance : AppliedModelingLib.FairDivision.AdditiveChoreInstance (Fin n) (Fin m)),
          AppliedModelingLib.FairDivision.IsOneOrRChoreCost choreInstance.cost r ∧
            ∀ (allocation : AppliedModelingLib.FairDivision.Allocation (Fin n) (Fin m)),
              choreInstance.IsFeasible allocation →
                choreInstance.IsEFX allocation → ¬choreInstance.IsParetoOptimal allocation
\end{ReviewVerbatim}



\paragraph{Lean proof endpoint (built separately)}\begin{ReviewVerbatim}
HT26EFXChores.efx_pareto_incompatibility
\end{ReviewVerbatim}

\paragraph{Recorded source-to-Spec screening}
\textbf{Verdict:} \texttt{matches}
\\
{\raggedright
\textbf{Reason:} Exact existence of a normalized one-\allowbreak{}or-\allowbreak{}r bi-\allowbreak{}valued feasible instance whose EFX allocations are not Pareto-\allowbreak{}optimal.\allowbreak{}
\par}
\reviewmatch{Matches source input}
\noindent\textbf{Reviewer annotation}\par
\reviewerbox
\clearpage
\hypertarget{source-claim-6906c411c024867a}{}
\section*{Review row 5}
\textbf{Source locator:} {\footnotesize\raggedright cited publication, lines 448-\allowbreak{}452\par}
\paragraph{Verbatim source input}
\begin{ReviewVerbatim}
Notice that an EFX allocation exists in this instance.
For each agent $i\in T_1$, she receives a large item from $A$ and a small item from $B$.
For all but one agent in $T_2$, each of them receives a large item from $A$ and a small item from $C$.
The remaining agent in $T_2$ receives one item from $B$ and two items from $C$.
The readers can easily verify this is an EFX allocation.

[Next verbatim source excerpt]

Agents are partitioned into two groups $T_1$ and $T_2$ with $|T_1|=t_1$ and $|T_2|=t_2$.
There are $m=2n+1$ items which are partitioned into three groups $A,B$, and $C$.
$A$ is the set of large items containing $n-1$ items which have cost $r$ to every agent.
$B$ contains $t_1+1$ items.
Each agent in $T_1$ values each item in $B$ at $1$ and each item in $C$ at $r$.
$C$ contains $t_2+1$ items.
Each agent in $T_2$ values each item in $B$ at $r$ and each item in $C$ at $1$.
Recall that $r$ can be any real number greater than $t_2+1$.
Table~\ref{tab:po} illustrates the instance.

\begin{table}[h]
    \centering
    \caption{A bi-valued instance where every EFX allocation fails to be Pareto-optimal.  Each entry is the cost of every item in the corresponding item class.}
    \begin{tabular}{cccc}
    \hline
         & $A=\{a_1,\ldots,a_{n-1}\}$ & $B=\{b_1,\ldots,b_{t_1+1}\}$ & $C=\{c_1,\ldots,c_{t_2+1}\}$\\
    \hline
       Agents in $T_1$ & $r$ & $1$ & $r$\\
       Agents in $T_2$ & $r$ & $r$ & $1$\\
    \hline
    \end{tabular}

    \label{tab:po}
\end{table}
\end{ReviewVerbatim}



\paragraph{Expanded PaperInterface specification}
\begin{ReviewVerbatim}
∀ (n : ℕ),
  4 ≤ n →
    ∀ (r : ℝ),
      ↑((n + 1) / 2) + 1 < r →
        ∃ (allocation : AppliedModelingLib.FairDivision.Allocation (Fin n) (Fin (2 * n + 1))),
          AppliedModelingLib.FairDivision.IsAllocationOf allocation Finset.univ ∧
            AppliedModelingLib.FairDivision.EFXForChores
              (AppliedModelingLib.FairDivision.additiveChoreCost fun (agent : Fin n) (item : Fin (2 * n + 1)) =>
                if ↑item < n - 1 then r
                else if ↑agent < n / 2 then if ↑item < n + n / 2 then 1 else r else if ↑item < n + n / 2 then r else 1)
              allocation
\end{ReviewVerbatim}



\paragraph{Lean proof endpoint (built separately)}\begin{ReviewVerbatim}
HT26EFXChores.efx_pareto_construction_has_efx
\end{ReviewVerbatim}

\paragraph{Recorded source-to-Spec screening}
\textbf{Verdict:} \texttt{matches}
\\
{\raggedright
\textbf{Reason:} Finite encoding has the stated group sizes, A/\allowbreak{}B/\allowbreak{}C costs, r threshold, and EFX existence.\allowbreak{}
\par}
\reviewmatch{Matches source input}
\noindent\textbf{Reviewer annotation}\par
\reviewerbox
\clearpage
\hypertarget{source-claim-5145ed0daed8afc6}{}
\section*{Review row 6}
\textbf{Source locator:} {\footnotesize\raggedright cited publication, lines 461-\allowbreak{}463\par}
\paragraph{Verbatim source input}
\begin{ReviewVerbatim}
\begin{proposition}\label{prop:PO-exist-large}
    For every $X_i$, there exists $g\in X_i$ with $c_i(g)=r$.
\end{proposition}

[Next verbatim source excerpt]

\begin{theorem}\label{thm:EFXPO}
    For every $n\geq 4$ and $r>\lceil n/2\rceil+1$, there exists a $(1,r)$-bi-valued instance with $n$ agents where every EFX allocation is not Pareto-optimal.
\end{theorem}

[Next verbatim source excerpt]

Let $t_1=\lfloor n/2\rfloor$ and $t_2=\lceil n/2\rceil$.
Agents are partitioned into two groups $T_1$ and $T_2$ with $|T_1|=t_1$ and $|T_2|=t_2$.
There are $m=2n+1$ items which are partitioned into three groups $A,B$, and $C$.
$A$ is the set of large items containing $n-1$ items which have cost $r$ to every agent.
$B$ contains $t_1+1$ items.
Each agent in $T_1$ values each item in $B$ at $1$ and each item in $C$ at $r$.
$C$ contains $t_2+1$ items.
Each agent in $T_2$ values each item in $B$ at $r$ and each item in $C$ at $1$.
Recall that $r$ can be any real number greater than $t_2+1$.
\end{ReviewVerbatim}



\paragraph{Expanded PaperInterface specification}
\begin{ReviewVerbatim}
∀ (n : ℕ),
  4 ≤ n →
    ∀ (r : ℝ),
      ↑((n + 1) / 2) + 1 < r →
        ∀ (allocation : AppliedModelingLib.FairDivision.Allocation (Fin n) (Fin (2 * n + 1))),
          AppliedModelingLib.FairDivision.IsAllocationOf allocation Finset.univ →
            AppliedModelingLib.FairDivision.EFXForChores
                (AppliedModelingLib.FairDivision.additiveChoreCost fun (agent : Fin n) (item : Fin (2 * n + 1)) =>
                  if ↑item < n - 1 then r
                  else
                    if ↑agent < n / 2 then if ↑item < n + n / 2 then 1 else r else if ↑item < n + n / 2 then r else 1)
                allocation →
              ∀ (agent : Fin n),
                ∃ item ∈ allocation agent,
                  (if ↑item < n - 1 then r
                    else
                      if ↑agent < n / 2 then if ↑item < n + n / 2 then 1 else r
                      else if ↑item < n + n / 2 then r else 1) =
                    r
\end{ReviewVerbatim}



\paragraph{Lean proof endpoint (built separately)}\begin{ReviewVerbatim}
HT26EFXChores.efx_po_every_agent_large
\end{ReviewVerbatim}

\paragraph{Recorded source-to-Spec screening}
\textbf{Verdict:} \texttt{matches}
\\
{\raggedright
\textbf{Reason:} Uses the stated instance and asserts every agent's bundle contains an item of own cost r.\allowbreak{}
\par}
\reviewmatch{Matches source input}
\noindent\textbf{Reviewer annotation}\par
\reviewerbox
\clearpage
\hypertarget{source-claim-8a3c68dd9df286a5}{}
\section*{Review row 7}
\textbf{Source locator:} {\footnotesize\raggedright cited publication, lines 502-\allowbreak{}504\par}
\paragraph{Verbatim source input}
\begin{ReviewVerbatim}
\begin{proposition}\label{prop:r+1r+2}
    For every agent $i$, we have $c_i(X_i)\geq r+1$. Moreover, there exists an agent $i$ with $c_i(X_i)\geq r+2$.
\end{proposition}

[Next verbatim source excerpt]

\begin{theorem}\label{thm:EFXPO}
    For every $n\geq 4$ and $r>\lceil n/2\rceil+1$, there exists a $(1,r)$-bi-valued instance with $n$ agents where every EFX allocation is not Pareto-optimal.
\end{theorem}

[Next verbatim source excerpt]

Let $t_1=\lfloor n/2\rfloor$ and $t_2=\lceil n/2\rceil$.
Agents are partitioned into two groups $T_1$ and $T_2$ with $|T_1|=t_1$ and $|T_2|=t_2$.
There are $m=2n+1$ items which are partitioned into three groups $A,B$, and $C$.
$A$ is the set of large items containing $n-1$ items which have cost $r$ to every agent.
$B$ contains $t_1+1$ items.
Each agent in $T_1$ values each item in $B$ at $1$ and each item in $C$ at $r$.
$C$ contains $t_2+1$ items.
Each agent in $T_2$ values each item in $B$ at $r$ and each item in $C$ at $1$.
Recall that $r$ can be any real number greater than $t_2+1$.
\end{ReviewVerbatim}



\paragraph{Expanded PaperInterface specification}
\begin{ReviewVerbatim}
∀ (n : ℕ),
  4 ≤ n →
    ∀ (r : ℝ),
      ↑((n + 1) / 2) + 1 < r →
        ∀ (allocation : AppliedModelingLib.FairDivision.Allocation (Fin n) (Fin (2 * n + 1))),
          AppliedModelingLib.FairDivision.IsAllocationOf allocation Finset.univ →
            AppliedModelingLib.FairDivision.EFXForChores
                (AppliedModelingLib.FairDivision.additiveChoreCost fun (agent : Fin n) (item : Fin (2 * n + 1)) =>
                  if ↑item < n - 1 then r
                  else
                    if ↑agent < n / 2 then if ↑item < n + n / 2 then 1 else r else if ↑item < n + n / 2 then r else 1)
                allocation →
              (∀ (agent : Fin n),
                  r + 1 ≤
                    AppliedModelingLib.FairDivision.additiveChoreCost
                      (fun (agent : Fin n) (item : Fin (2 * n + 1)) =>
                        if ↑item < n - 1 then r
                        else
                          if ↑agent < n / 2 then if ↑item < n + n / 2 then 1 else r
                          else if ↑item < n + n / 2 then r else 1)
                      agent (allocation agent)) ∧
                ∃ (agent : Fin n),
                  r + 2 ≤
                    AppliedModelingLib.FairDivision.additiveChoreCost
                      (fun (agent : Fin n) (item : Fin (2 * n + 1)) =>
                        if ↑item < n - 1 then r
                        else
                          if ↑agent < n / 2 then if ↑item < n + n / 2 then 1 else r
                          else if ↑item < n + n / 2 then r else 1)
                      agent (allocation agent)
\end{ReviewVerbatim}



\paragraph{Lean proof endpoint (built separately)}\begin{ReviewVerbatim}
HT26EFXChores.efx_po_cost_lower_bounds
\end{ReviewVerbatim}

\paragraph{Recorded source-to-Spec screening}
\textbf{Verdict:} \texttt{matches}
\\
{\raggedright
\textbf{Reason:} Uses the specified EFXPO instance and concludes every own cost is at least r+1, with one at least r+2.\allowbreak{}
\par}
\reviewmatch{Matches source input}
\noindent\textbf{Reviewer annotation}\par
\reviewerbox
\clearpage
\hypertarget{source-claim-03ed4c5388a9d78f}{}
\section*{Review row 8}
\textbf{Source locator:} {\footnotesize\raggedright cited publication, lines 532-\allowbreak{}534\par}
\paragraph{Verbatim source input}
\begin{ReviewVerbatim}
\begin{theorem}\label{thm:four}
    For four agents, there exists an EFX allocation in every bi-valued instance.
\end{theorem}

[Next verbatim source excerpt]

\section{Preliminaries}
A set $M$ of $m$ chores $g_1,\ldots,g_m$ is allocated to a set $N$ of $n$ agents $1,\ldots,n$.
We will use ``item'' and ``chore'' interchangeably in this paper to refer to an element in $M$.
Each agent $i$'s cost function $c_i:2^M\to\mathbb{R}_{\geq0}$ is additive:
$$c_i(S)=\sum_{g\in S}c_i(\{g\}).$$
We write $c_i(g_j)$ in short for $c_i(\{g_j\})$ for notation simplicity.
We will use ``cost function'' and ``disutility function'' interchangeably in this paper.

[Next verbatim source excerpt]

We say that the cost function is tri-valued if there exist constants $p,q,r\geq 0$ such that $c_i(g_j)\in\{p,q,r\}$ for any $i\in N$ and $g_j\in M$.
The cost function is bi-valued if there exist constants $p,q\geq 0$ such that $c_i(g_j)\in\{p,q\}$ for any $i\in N$ and $g_j\in M$.
For bi-valued cost functions, we will focus on the case $0<p<q$ (the case $0=p<q$ corresponds to the setting with binary cost functions and we know an allocation that is EFX and Pareto-optimal always exists), and we will assume $c_i(g_j)\in\{1,r\}$ (for any $i\in N$ and $g_j\in M$) instead.
We say that the item $g_j$ is ``large'' for agent $i$ if $c_i(g_j)=r$ and it is ``small'' for agent $i$ if $c_i(g_j)=1$.
\end{ReviewVerbatim}



\paragraph{Expanded PaperInterface specification}
\begin{ReviewVerbatim}
∀ (Item : Type) (choreInstance : AppliedModelingLib.FairDivision.AdditiveChoreInstance (Fin 4) Item),
  AppliedModelingLib.FairDivision.IsBiValuedChoreCost choreInstance.cost →
    ∃ (allocation : AppliedModelingLib.FairDivision.Allocation (Fin 4) Item),
      choreInstance.IsFeasible allocation ∧ choreInstance.IsEFX allocation
\end{ReviewVerbatim}



\paragraph{Lean proof endpoint (built separately)}\begin{ReviewVerbatim}
HT26EFXChores.four_agent_bi_valued_efx_exists
\end{ReviewVerbatim}

\paragraph{Recorded source-to-Spec screening}
\textbf{Verdict:} \texttt{matches}
\\
{\raggedright
\textbf{Reason:} Exact four-\allowbreak{}agent bi-\allowbreak{}valued EFX-\allowbreak{}existence claim;\allowbreak{} instance nonnegativity is carried by AdditiveChoreInstance.\allowbreak{}
\par}
\reviewmatch{Matches source input}
\noindent\textbf{Reviewer annotation}\par
\reviewerbox
\clearpage
\hypertarget{source-claim-e126eba4635c1cd5}{}
\section*{Review row 9}
\textbf{Source locator:} {\footnotesize\raggedright cited publication, lines 577-\allowbreak{}581\par}
\paragraph{Verbatim source input}
\begin{ReviewVerbatim}
\begin{lemma}[Insertion of an $M_{34}$ item]\label{lem:insertion}
Let $X$ be an EFX allocation of a set of chores $T$.  Let $g\in M_{34}$ be a
chore that is small for at least three agents.  Then there is an EFX
allocation of $T\cup\{g\}$.
\end{lemma}

[Next verbatim source excerpt]

We say that the cost function is tri-valued if there exist constants $p,q,r\geq 0$ such that $c_i(g_j)\in\{p,q,r\}$ for any $i\in N$ and $g_j\in M$.
The cost function is bi-valued if there exist constants $p,q\geq 0$ such that $c_i(g_j)\in\{p,q\}$ for any $i\in N$ and $g_j\in M$.
For bi-valued cost functions, we will focus on the case $0<p<q$ (the case $0=p<q$ corresponds to the setting with binary cost functions and we know an allocation that is EFX and Pareto-optimal always exists), and we will assume $c_i(g_j)\in\{1,r\}$ (for any $i\in N$ and $g_j\in M$) instead.
We say that the item $g_j$ is ``large'' for agent $i$ if $c_i(g_j)=r$ and it is ``small'' for agent $i$ if $c_i(g_j)=1$.

[Next verbatim source excerpt]

We will assume $r>2$ in the remaining part of this section.
For a chore $g$, let
\[
 S(g)=\{i\in N:c_i(g)=1\}
\]
be the set of agents for whom $g$ is small.  We partition the chores into
\[
M_{01}=\{g:|S(g)|\le 1\},\qquad
M_2=\{g:|S(g)|=2\},\qquad
M_{34}=\{g:|S(g)|\ge 3\}.
\]
For an item in $M_2$, we write $(i,j)$ for an item small exactly for agents
$i$ and $j$.  Thus the items in $M_2$ form a multigraph on the vertex set
$N$, where an edge $(i,j)$ is a chore small for precisely agents $i$ and $j$.
\end{ReviewVerbatim}



\paragraph{Expanded PaperInterface specification}
\begin{ReviewVerbatim}
∀ (Item : Type) (r : ℝ) (cost : AppliedModelingLib.FairDivision.ChoreCost (Fin 4) Item) (chores : Finset Item)
  (item : Item),
  2 < r →
    AppliedModelingLib.FairDivision.IsOneOrRChoreCost cost r →
      item ∉ chores →
        AppliedModelingLib.FairDivision.IsSmallForAtLeastThree cost item →
          ∀ (allocation : AppliedModelingLib.FairDivision.Allocation (Fin 4) Item),
            AppliedModelingLib.FairDivision.IsAllocationOf allocation chores →
              AppliedModelingLib.FairDivision.EFXForChores (AppliedModelingLib.FairDivision.additiveChoreCost cost)
                  allocation →
                ∃ (extended : AppliedModelingLib.FairDivision.Allocation (Fin 4) Item),
                  AppliedModelingLib.FairDivision.IsAllocationOf extended (insert item chores) ∧
                    AppliedModelingLib.FairDivision.EFXForChores
                      (AppliedModelingLib.FairDivision.additiveChoreCost cost) extended
\end{ReviewVerbatim}



\paragraph{Lean proof endpoint (built separately)}\begin{ReviewVerbatim}
HT26EFXChores.m34_insertion
\end{ReviewVerbatim}

\paragraph{Recorded source-to-Spec screening}
\textbf{Verdict:} \texttt{matches}
\\
{\raggedright
\textbf{Reason:} Exact four-\allowbreak{}agent M34 insertion result;\allowbreak{} the explicit freshness condition is standard nonsemantic insertion presentation.\allowbreak{}
\par}
\reviewmatch{Matches source input}
\noindent\textbf{Reviewer annotation}\par
\reviewerbox
\clearpage
\hypertarget{source-claim-5fbcb5b015031605}{}
\section*{Review row 10}
\textbf{Source locator:} {\footnotesize\raggedright cited publication, lines 673-\allowbreak{}677\par}
\paragraph{Verbatim source input}
\begin{ReviewVerbatim}
\begin{remark}
    This idea also works for general numbers of agents.
    If we have a partial EFX allocation such that each of the remaining unallocated items is valued as large for at most one agent, then we can always allocate the remaining items one by one until we get a complete EFX allocation.
    Therefore, under bi-valued setting, items valued as large by at most one agent are never a real obstacle for EFX allocations.
\end{remark}

[Next verbatim source excerpt]

We say that the cost function is tri-valued if there exist constants $p,q,r\geq 0$ such that $c_i(g_j)\in\{p,q,r\}$ for any $i\in N$ and $g_j\in M$.
The cost function is bi-valued if there exist constants $p,q\geq 0$ such that $c_i(g_j)\in\{p,q\}$ for any $i\in N$ and $g_j\in M$.
For bi-valued cost functions, we will focus on the case $0<p<q$ (the case $0=p<q$ corresponds to the setting with binary cost functions and we know an allocation that is EFX and Pareto-optimal always exists), and we will assume $c_i(g_j)\in\{1,r\}$ (for any $i\in N$ and $g_j\in M$) instead.
We say that the item $g_j$ is ``large'' for agent $i$ if $c_i(g_j)=r$ and it is ``small'' for agent $i$ if $c_i(g_j)=1$.

[Next verbatim source excerpt]

\section{EFX Allocations for Four Agents with Bi-Valued Disutilities}
We continue to study bi-valued instances.
In this section, we show that an EFX allocation always exists for four agents.

\begin{theorem}\label{thm:four}
    For four agents, there exists an EFX allocation in every bi-valued instance.
\end{theorem}

For $r=2$, \citet{lin2025approximately} shows that EFX is compatible with Pareto-optimality for any number of agents.
If we only consider EFX, the case with $r\leq 2$ can be handled by the following simple variant of the round-robin algorithm for any number of agents.
Let $m=an+b$ for $b<n$.
If $b=0$, the standard round-robin algorithm outputs an EFX allocation $X=(X_1,\ldots,X_n)$: for any pair of agents $i,j$ with $i<j$, $i$ will not envy $j$ as $i$ always picks first; we also have $c_j(X_j)-c_j(X_i)\leq r-1$ (if in a round agent $j$ begins to receive an large item whereas agent $i$'s item in this round is small for agent $j$, then every item in the remaining process of this algorithm is large for agent $j$), $j$ will not envy $i$ up to removing any item from $j$'s bundle since $r-1\leq1$.
If $b>0$, we introduce a special ``first reverse round'' of the round-robin algorithm where agents $b,b-1,\ldots,1$ iteratively pick favorite items in the order.
After that, we perform the standard round-robin algorithm in the standard order $1,\ldots,n$.
For each pair of agents $i,j$ with $i<j$, agent $i$ will not envy agent $j$ after removing the item picked in the first special round. If this item is a small item for agent $i$, EFX condition holds; if this item is a large item, then $X_i$ only contains large items, and EFX condition also holds.
The reason agent $j$ will not strongly envy agent $i$ is similar as the $b=0$ case: agent $j$ has an advantage to agent $i$ in the first round, and after that the disadvantage is at most $r-1$.

We will assume $r>2$ in the remaining part of this section.
For a chore $g$, let
\[
 S(g)=\{i\in N:c_i(g)=1\}
\]
be the set of agents for whom $g$ is small.  We partition the chores into
\[
M_{01}=\{g:|S(g)|\le 1\},\qquad
M_2=\{g:|S(g)|=2\},\qquad
M_{34}=\{g:|S(g)|\ge 3\}.
\]
For an item in $M_2$, we write $(i,j)$ for an item small exactly for agents
\end{ReviewVerbatim}



\paragraph{Expanded PaperInterface specification}
\begin{ReviewVerbatim}
∀ (Agent Item : Type) [Fintype Agent] [Nonempty Agent] [DecidableEq Agent] [inst : DecidableEq Item] (r : ℝ)
  (cost : AppliedModelingLib.FairDivision.ChoreCost Agent Item) (remaining chores : Finset Item),
  2 < r →
    AppliedModelingLib.FairDivision.IsOneOrRChoreCost cost r →
      Disjoint remaining chores →
        (∀ item ∈ remaining, ∀ (first second : Agent), cost first item = r → cost second item = r → first = second) →
          (∀ item ∈ remaining,
              ∀ (allocation : AppliedModelingLib.FairDivision.Allocation Agent Item),
                AppliedModelingLib.FairDivision.IsAllocationOf allocation chores →
                  AppliedModelingLib.FairDivision.EFXForChores (AppliedModelingLib.FairDivision.additiveChoreCost cost)
                      allocation →
                    ∃ (extended : AppliedModelingLib.FairDivision.Allocation Agent Item),
                      AppliedModelingLib.FairDivision.IsAllocationOf extended (insert item chores) ∧
                        AppliedModelingLib.FairDivision.EFXForChores
                          (AppliedModelingLib.FairDivision.additiveChoreCost cost) extended) ∧
            ∀ (allocation : AppliedModelingLib.FairDivision.Allocation Agent Item),
              AppliedModelingLib.FairDivision.IsAllocationOf allocation chores →
                AppliedModelingLib.FairDivision.EFXForChores (AppliedModelingLib.FairDivision.additiveChoreCost cost)
                    allocation →
                  ∃ (extended : AppliedModelingLib.FairDivision.Allocation Agent Item),
                    AppliedModelingLib.FairDivision.IsAllocationOf extended (chores ∪ remaining) ∧
                      AppliedModelingLib.FairDivision.EFXForChores
                        (AppliedModelingLib.FairDivision.additiveChoreCost cost) extended
\end{ReviewVerbatim}



\paragraph{Lean proof endpoint (built separately)}\begin{ReviewVerbatim}
HT26EFXChores.m34_insertion_general_agents
\end{ReviewVerbatim}

\paragraph{Recorded source-to-Spec screening}
\textbf{Verdict:} \texttt{matches}
\\
{\raggedright
\textbf{Reason:} Exact general finite-\allowbreak{}agent sequential insertion claim for remaining items large for at most one agent, including final EFX completion.\allowbreak{}
\par}
\reviewmatch{Matches source input}
\noindent\textbf{Reviewer annotation}\par
\reviewerbox
\clearpage
\hypertarget{source-claim-58fb9c46acb21eca}{}
\section*{Review row 11}
\textbf{Source locator:} {\footnotesize\raggedright cited publication, lines 706-\allowbreak{}724\par}
\paragraph{Verbatim source input}
\begin{ReviewVerbatim}
\begin{lemma}[Composition]\label{lem:composition}
Let $P=(P_1,\ldots,P_n)$ be an allocation of some items and $B=(B_1,\ldots,B_n)$ an allocation of disjoint
items.  If
\[
 c_i(P_i)\le c_i(P_j)\qquad\forall i,j
\]
and
\[
 c_i(B_i)-1\le c_i(B_j)\qquad\forall i,j,
\]
then the combined allocation $X_i=P_i\cup B_i$ is EFX. %provided every agent whose second inequality is tight has a small item in her final bundle.
More generally, $X$ is EFX whenever
\[
 c_i(P_i)-c_i(P_j)
 \le c_i(B_j)-c_i(B_i)+\min_{g\in X_i}c_i(g)
 \qquad\forall i,j.
\]
If the inequalities above hold only for a particular pair of agents $(i,j)$, then $i$ does not strongly envy $j$.
\end{lemma}

[Next verbatim source excerpt]

We say that the cost function is tri-valued if there exist constants $p,q,r\geq 0$ such that $c_i(g_j)\in\{p,q,r\}$ for any $i\in N$ and $g_j\in M$.
The cost function is bi-valued if there exist constants $p,q\geq 0$ such that $c_i(g_j)\in\{p,q\}$ for any $i\in N$ and $g_j\in M$.
For bi-valued cost functions, we will focus on the case $0<p<q$ (the case $0=p<q$ corresponds to the setting with binary cost functions and we know an allocation that is EFX and Pareto-optimal always exists), and we will assume $c_i(g_j)\in\{1,r\}$ (for any $i\in N$ and $g_j\in M$) instead.
We say that the item $g_j$ is ``large'' for agent $i$ if $c_i(g_j)=r$ and it is ``small'' for agent $i$ if $c_i(g_j)=1$.

[Next verbatim source excerpt]

We will assume $r>2$ in the remaining part of this section.
\end{ReviewVerbatim}



\paragraph{Expanded PaperInterface specification}
\begin{ReviewVerbatim}
∀ (Agent Item : Type) (r : ℝ) (cost : AppliedModelingLib.FairDivision.ChoreCost Agent Item)
  (leftChores rightChores : Finset Item)
  (leftAllocation rightAllocation : AppliedModelingLib.FairDivision.Allocation Agent Item),
  2 < r →
    AppliedModelingLib.FairDivision.IsOneOrRChoreCost cost r →
      Disjoint leftChores rightChores →
        AppliedModelingLib.FairDivision.IsAllocationOf leftAllocation leftChores →
          AppliedModelingLib.FairDivision.IsAllocationOf rightAllocation rightChores →
            AppliedModelingLib.FairDivision.IsAllocationOf
                (fun (agent : Agent) => leftAllocation agent ∪ rightAllocation agent) (leftChores ∪ rightChores) ∧
              ((AppliedModelingLib.FairDivision.EnvyFreeForChores
                      (AppliedModelingLib.FairDivision.additiveChoreCost cost) leftAllocation ∧
                    ∀ (i j : Agent),
                      AppliedModelingLib.FairDivision.additiveChoreCost cost i (rightAllocation i) - 1 ≤
                        AppliedModelingLib.FairDivision.additiveChoreCost cost i (rightAllocation j)) →
                  AppliedModelingLib.FairDivision.EFXForChores (AppliedModelingLib.FairDivision.additiveChoreCost cost)
                    fun (agent : Agent) => leftAllocation agent ∪ rightAllocation agent) ∧
                ((∀ (i j : Agent) (hnonempty : Finset.Nonempty (leftAllocation i ∪ rightAllocation i)),
                      AppliedModelingLib.FairDivision.additiveChoreCost cost i (leftAllocation i) -
                          AppliedModelingLib.FairDivision.additiveChoreCost cost i (leftAllocation j) ≤
                        AppliedModelingLib.FairDivision.additiveChoreCost cost i (rightAllocation j) -
                            AppliedModelingLib.FairDivision.additiveChoreCost cost i (rightAllocation i) +
                          (Finset.image (cost i) (leftAllocation i ∪ rightAllocation i)).min'
                            (HT26EFXChores.compositionSpec._proof_1 Agent Item cost leftAllocation rightAllocation i
                              hnonempty)) →
                    AppliedModelingLib.FairDivision.EFXForChores
                      (AppliedModelingLib.FairDivision.additiveChoreCost cost) fun (agent : Agent) =>
                      leftAllocation agent ∪ rightAllocation agent) ∧
                  ∀ (i j : Agent) (hnonempty : Finset.Nonempty (leftAllocation i ∪ rightAllocation i)),
                    AppliedModelingLib.FairDivision.additiveChoreCost cost i (leftAllocation i) -
                          AppliedModelingLib.FairDivision.additiveChoreCost cost i (leftAllocation j) ≤
                        AppliedModelingLib.FairDivision.additiveChoreCost cost i (rightAllocation j) -
                            AppliedModelingLib.FairDivision.additiveChoreCost cost i (rightAllocation i) +
                          (Finset.image (cost i) (leftAllocation i ∪ rightAllocation i)).min'
                            (HT26EFXChores.compositionSpec._proof_1 Agent Item cost leftAllocation rightAllocation i
                              hnonempty) →
                      AppliedModelingLib.FairDivision.DoesNotStronglyEnvyForChores
                        (AppliedModelingLib.FairDivision.additiveChoreCost cost)
                        (fun (agent : Agent) => leftAllocation agent ∪ rightAllocation agent) i j
\end{ReviewVerbatim}



\paragraph{Lean proof endpoint (built separately)}\begin{ReviewVerbatim}
HT26EFXChores.composition
\end{ReviewVerbatim}

\paragraph{Recorded source-to-Spec screening}
\textbf{Verdict:} \texttt{matches}
\\
{\raggedright
\textbf{Reason:} Exact disjoint-\allowbreak{}allocation composition, envy-\allowbreak{}free-\allowbreak{}plus-\allowbreak{}minus-\allowbreak{}one clause, minimum-\allowbreak{}item clause, and no-\allowbreak{}strong-\allowbreak{}envy consequence.\allowbreak{}
\par}
\reviewmatch{Matches source input}
\noindent\textbf{Reviewer annotation}\par
\reviewerbox
\clearpage
\hypertarget{source-claim-1b77c29da35b62c2}{}
\section*{Review row 12}
\textbf{Source locator:} {\footnotesize\raggedright cited publication, lines 760-\allowbreak{}767\par}
\paragraph{Verbatim source input}
\begin{ReviewVerbatim}
\begin{lemma}\label{lem:canonical}
    We have the following properties regarding canonical allocations.
    \begin{enumerate}[label=(\alph*)]
        \item Every canonical allocation is EFX. For $b=0$, every canonical allocation is envy-free.
        \item For any $b>0$, there always exists a super-canonical allocation.
        \item If $a>0$ and $b=2$, for any partition $N_1,N_2$ of $N$ with $|N_1|=|N_2|=2$, there exist $i\in N_1$ and $j\in N_2$ such that there exists a super-canonical allocation $P$ with $|P_i|=|P_j|=a$.
    \end{enumerate}
\end{lemma}

[Next verbatim source excerpt]

We will assume $r>2$ in the remaining part of this section.
For a chore $g$, let
\[
 S(g)=\{i\in N:c_i(g)=1\}
\]
be the set of agents for whom $g$ is small.  We partition the chores into
\[
M_{01}=\{g:|S(g)|\le 1\},\qquad
M_2=\{g:|S(g)|=2\},\qquad
M_{34}=\{g:|S(g)|\ge 3\}.
\]
For an item in $M_2$, we write $(i,j)$ for an item small exactly for agents
$i$ and $j$.  Thus the items in $M_2$ form a multigraph on the vertex set
$N$, where an edge $(i,j)$ is a chore small for precisely agents $i$ and $j$.

[Next verbatim source excerpt]

For each agent $i$, let
\[
 n_i=\left|\{g\in M_{01}:S(g)=\{i\}\}\right|
\]
be the number of $M_{01}$ items uniquely small for $i$.  Items in $M_{01}$
with $S(g)=\emptyset$ are all-large items.
Let $|M_{01}|=4a+b$ with $b\in\{0,1,2,3\}$.
As mentioned before, we will make sure each agent receives either $a$ items or $a+1$ items.

Fix quotas $q_i\in\{a,a+1\}$ with $\sum_i q_i=|M_{01}|$.
We consider allocations of $M_{01}$ with quotas $q=(q_1,q_2,q_3,q_4)$ in which each agent $i$ first receives as many of her uniquely-small items as possible, up to quota $q_i$, and the remaining slots are filled with the remaining items.
We will show that all such allocations satisfy EFX.
It is also important that we have the freedom to choose which agents receive $a$ items and which agents receive $a+1$ items.
This freedom enables us to append the allocation of $M_2$ after the allocation of $M_{01}$.

We begin by defining this type of allocations.


[Next verbatim source excerpt]

be the set of agents for whom $g$ is small.  We partition the chores into
\[
M_{01}=\{g:|S(g)|\le 1\},\qquad
M_2=\{g:|S(g)|=2\},\qquad
M_{34}=\{g:|S(g)|\ge 3\}.
\]
For an item in $M_2$, we write $(i,j)$ for an item small exactly for agents
$i$ and $j$.  Thus the items in $M_2$ form a multigraph on the vertex set
$N$, where an edge $(i,j)$ is a chore small for precisely agents $i$ and $j$.

[Next verbatim source excerpt]

\begin{definition}
    An allocation $P=(P_1,P_2,P_3,P_4)$ of $M_{01}$ is \emph{canonical} if, for each agent $i$, $|P_i|\in\{a,a+1\}$ and $P_i$ contains $\min\{|P_i|, n_i\}$ items that are small for agent $i$.
    In a canonical allocation with $b>0$, agents receiving $a$ items are called \emph{short} agents, and agents receiving $a+1$ items are called \emph{long} agents.
    For $b=0$ and $a>0$, every agent in a canonical is both long and short.
    For $b>0$, an allocation $P=(P_1,P_2,P_3,P_4)$ is \emph{super-canonical} if it is canonical and satisfies $c_i(P_i)\leq c_i(P_j)-r$ for every short agent $i$ and every long agent $j$.
\end{definition}
\end{ReviewVerbatim}



\paragraph{Expanded PaperInterface specification}
\begin{ReviewVerbatim}
∀ (Item : Type) (r : ℝ) (cost : AppliedModelingLib.FairDivision.ChoreCost (Fin 4) Item) (chores : Finset Item)
  (a b : ℕ),
  2 < r →
    AppliedModelingLib.FairDivision.IsOneOrRChoreCost cost r →
      b ≤ 3 →
        chores.card = 4 * a + b →
          (∀ item ∈ chores, AppliedModelingLib.FairDivision.IsSmallForAtMostOne cost item) →
            (∀ (quota : Fin 4 → ℕ) (allocation : AppliedModelingLib.FairDivision.Allocation (Fin 4) Item),
                (∀ (agent : Fin 4), quota agent = a ∨ quota agent = a + 1) →
                  AppliedModelingLib.FairDivision.IsCanonicalSmallChoreAllocation cost chores quota allocation →
                    AppliedModelingLib.FairDivision.EFXForChores
                      (AppliedModelingLib.FairDivision.additiveChoreCost cost) allocation) ∧
              (b = 0 →
                  ∀ (quota : Fin 4 → ℕ) (allocation : AppliedModelingLib.FairDivision.Allocation (Fin 4) Item),
                    (∀ (agent : Fin 4), quota agent = a) →
                      AppliedModelingLib.FairDivision.IsCanonicalSmallChoreAllocation cost chores quota allocation →
                        AppliedModelingLib.FairDivision.EnvyFreeForChores
                          (AppliedModelingLib.FairDivision.additiveChoreCost cost) allocation) ∧
                (0 < b →
                    ∃ (quota : Fin 4 → ℕ) (allocation : AppliedModelingLib.FairDivision.Allocation (Fin 4) Item),
                      (∀ (agent : Fin 4), quota agent = a ∨ quota agent = a + 1) ∧
                        AppliedModelingLib.FairDivision.IsCanonicalSmallChoreAllocation cost chores quota allocation ∧
                          ∀ (i j : Fin 4),
                            quota i = a →
                              quota j = a + 1 →
                                AppliedModelingLib.FairDivision.additiveChoreCost cost i (allocation i) ≤
                                  AppliedModelingLib.FairDivision.additiveChoreCost cost i (allocation j) - r) ∧
                  (0 < a →
                    b = 2 →
                      ∀ (N1 N2 : Finset (Fin 4)),
                        N1 ∪ N2 = Finset.univ →
                          Disjoint N1 N2 →
                            N1.card = 2 →
                              N2.card = 2 →
                                ∃ i ∈ N1,
                                  ∃ j ∈ N2,
                                    ∃ (quota : Fin 4 → ℕ) (allocation :
                                      AppliedModelingLib.FairDivision.Allocation (Fin 4) Item),
                                      (∀ (agent : Fin 4), quota agent = a ∨ quota agent = a + 1) ∧
                                        AppliedModelingLib.FairDivision.IsCanonicalSmallChoreAllocation cost chores
                                            quota allocation ∧
                                          (∀ (x y : Fin 4),
                                              quota x = a →
                                                quota y = a + 1 →
                                                  AppliedModelingLib.FairDivision.additiveChoreCost cost x
                                                      (allocation x) ≤
                                                    AppliedModelingLib.FairDivision.additiveChoreCost cost x
                                                        (allocation y) -
                                                      r) ∧
                                            quota i = a ∧ quota j = a)
\end{ReviewVerbatim}



\paragraph{Lean proof endpoint (built separately)}\begin{ReviewVerbatim}
HT26EFXChores.canonical_allocation_properties
\end{ReviewVerbatim}

\paragraph{Recorded source-to-Spec screening}
\textbf{Verdict:} \texttt{matches}
\\
{\raggedright
\textbf{Reason:} Captures all canonical and super-\allowbreak{}canonical clauses;\allowbreak{} feasibility and quota-\allowbreak{}cardinality are within the expanded canonical predicate.\allowbreak{}
\par}
\reviewmatch{Matches source input}
\noindent\textbf{Reviewer annotation}\par
\reviewerbox
\clearpage
\hypertarget{source-claim-0869bc0832b92b55}{}
\section*{Review row 13}
\textbf{Source locator:} {\footnotesize\raggedright cited publication, lines 1019-\allowbreak{}1031\par}
\paragraph{Verbatim source input}
\begin{ReviewVerbatim}
\begin{lemma}[Balanced orientation]\label{lem:balancedorientation}
Let $G=(V,E)$ be a multigraph and let $(b_i)_{i\in V}$ be nonnegative
integers satisfying $\sum_{i\in V}b_i=|E|$.  There exists an assignment of
every edge to one of its endpoints such that vertex $i$ receives exactly
$b_i$ edges if and only if
\begin{equation}\label{eqn:orientation}
    \sum_{i\in U}b_i\le \operatorname{inc}_G(U)
 \qquad\text{for every }U\subseteq V.
\end{equation}

If we have \(\sum_{i\in V}b_i\le |E|\) instead, and \eqref{eqn:orientation} is satisfied, then one can assign some edges to their endpoints so that vertex \(i\)
receives exactly \(b_i\) edges, leaving the remaining edges unassigned.
\end{lemma}
\end{ReviewVerbatim}



\paragraph{Expanded PaperInterface specification}
\begin{ReviewVerbatim}
∀ (Vertex Edge : Type) [inst : Fintype Vertex] (edges : Finset Edge) (endpoints : Edge → Finset Vertex)
  (quota : Vertex → ℕ),
  (∀ edge ∈ edges, (endpoints edge).card = 2) →
    (Finset.univ.sum quota = edges.card →
        ((∃ (owner : Edge → Option Vertex),
            (∀ edge ∈ edges, ∃ (vertex : Vertex), owner edge = some vertex ∧ vertex ∈ endpoints edge) ∧
              ∀ (vertex : Vertex), {edge ∈ edges | owner edge = some vertex}.card = quota vertex) ↔
          ∀ (vertices : Finset Vertex),
            vertices.sum quota ≤ {edge ∈ edges | (endpoints edge ∩ vertices).Nonempty}.card)) ∧
      (Finset.univ.sum quota ≤ edges.card →
        (∀ (vertices : Finset Vertex),
            vertices.sum quota ≤ {edge ∈ edges | (endpoints edge ∩ vertices).Nonempty}.card) →
          ∃ (owner : Edge → Option Vertex),
            (∀ edge ∈ edges, ∀ (vertex : Vertex), owner edge = some vertex → vertex ∈ endpoints edge) ∧
              ∀ (vertex : Vertex), {edge ∈ edges | owner edge = some vertex}.card = quota vertex)
\end{ReviewVerbatim}



\paragraph{Lean proof endpoint (built separately)}\begin{ReviewVerbatim}
HT26EFXChores.balanced_orientation
\end{ReviewVerbatim}

\paragraph{Recorded source-to-Spec screening}
\textbf{Verdict:} \texttt{matches}
\\
{\raggedright
\textbf{Reason:} The expanded Spec gives the exact all-\allowbreak{}edges endpoint-\allowbreak{}assignment equivalence when total demand equals the edge count and the partial endpoint-\allowbreak{}assignment conclusion when demand is at most the edge count, with the same incident-\allowbreak{}edge inequalities for every vertex subset.\allowbreak{}
\par}
\reviewmatch{Matches source input}
\noindent\textbf{Reviewer annotation}\par
\reviewerbox
\clearpage
\hypertarget{source-claim-6263bf063357fb29}{}
\section*{Review row 14}
\textbf{Source locator:} {\footnotesize\raggedright cited publication, lines 1045-\allowbreak{}1047\par}
\paragraph{Verbatim source input}
\begin{ReviewVerbatim}
\begin{lemma}[EFX Allocation of $M_2$]\label{lem:M2}
There is an EFX allocation of $M_2$.
\end{lemma}

[Next verbatim source excerpt]

We will assume $r>2$ in the remaining part of this section.
For a chore $g$, let
\[
 S(g)=\{i\in N:c_i(g)=1\}
\]
be the set of agents for whom $g$ is small.  We partition the chores into
\[
M_{01}=\{g:|S(g)|\le 1\},\qquad
M_2=\{g:|S(g)|=2\},\qquad
M_{34}=\{g:|S(g)|\ge 3\}.
\]
For an item in $M_2$, we write $(i,j)$ for an item small exactly for agents
$i$ and $j$.  Thus the items in $M_2$ form a multigraph on the vertex set
$N$, where an edge $(i,j)$ is a chore small for precisely agents $i$ and $j$.

[Next verbatim source excerpt]

\subsection{The $M_2$ Multigraph Algorithm}
\label{subsec:2}
We now give a complete allocation procedure for the items in $M_2$.
Each item is small for some agents $i$ and $j$ and
large for the other two agents, and we say that the type of this item is $(i,j)$, or $ij$ for short.
In this section, we consider the case where the instance only has $M_{2}$ items.
We represent the instance by a
multigraph $R$ on $N=\{1,2,3,4\}$, where an item of type $(i,j)$ is an edge $(i,j)$.
We slightly abuse the notation by letting $R$ also represent the edge set, i.e., $R$ exactly corresponds to items in $M_2$.
Let $x_{ij}$ denote the multiplicity of edge $(i,j)$ and let $m_R=|R|$.

For $U\subseteq N$, define
\[
 \operatorname{inc}_R(U)=\left|\{e\in R:e\cap U\neq\emptyset\}\right|
\]
and
\[
 \Delta_R(U)=\frac{|U|m_R}{4}-\operatorname{inc}_R(U).
\]
We call $U$ \emph{deficient} if $\Delta_R(U)>0$.
\end{ReviewVerbatim}



\paragraph{Expanded PaperInterface specification}
\begin{ReviewVerbatim}
∀ (Item : Type) (r : ℝ) (cost : AppliedModelingLib.FairDivision.ChoreCost (Fin 4) Item) (chores : Finset Item),
  2 < r →
    AppliedModelingLib.FairDivision.IsOneOrRChoreCost cost r →
      (∀ item ∈ chores, AppliedModelingLib.FairDivision.IsSmallForExactlyTwo cost item) →
        ∃ (allocation : AppliedModelingLib.FairDivision.Allocation (Fin 4) Item),
          AppliedModelingLib.FairDivision.IsAllocationOf allocation chores ∧
            AppliedModelingLib.FairDivision.EFXForChores (AppliedModelingLib.FairDivision.additiveChoreCost cost)
              allocation
\end{ReviewVerbatim}



\paragraph{Lean proof endpoint (built separately)}\begin{ReviewVerbatim}
HT26EFXChores.m2_efx_allocation
\end{ReviewVerbatim}

\paragraph{Recorded source-to-Spec screening}
\textbf{Verdict:} \texttt{matches}
\\
{\raggedright
\textbf{Reason:} Exact four-\allowbreak{}agent one-\allowbreak{}or-\allowbreak{}r, r-\allowbreak{}greater-\allowbreak{}than-\allowbreak{}two, exactly-\allowbreak{}two-\allowbreak{}small-\allowbreak{}agent M2 premise and EFX allocation conclusion.\allowbreak{}
\par}
\reviewmatch{Matches source input}
\noindent\textbf{Reviewer annotation}\par
\reviewerbox
\clearpage
\hypertarget{source-claim-687ed4b91a799ce3}{}
\section*{Review row 15}
\textbf{Source locator:} {\footnotesize\raggedright cited publication, lines 1419-\allowbreak{}1462\par}
\paragraph{Verbatim source input}
\begin{ReviewVerbatim}
\begin{lemma}[Properties of the EFX allocations in Lemma~\ref{lem:M2}]\label{lem:M2properties}
The EFX allocation $B=(B_1,\ldots,B_4)$ ensured by Lemma~\ref{lem:M2} can be chosen to satisfy the
following additional properties.
\begin{enumerate}[label=(\roman*)]
\item Except in the two configurations described in~(ii), every possible envy is certified by removing a small item.  More precisely,
whenever $B_i$ contains an item that is small for agent $i$,
\[
 c_i(B_i)-1\le c_i(B_j)\qquad\forall j\in N,
\]
whereas every agent whose bundle contains only large items is envy-free:
\[
 c_i(B_i)\le c_i(B_j)\qquad\forall j\in N.
\]

\item Up to relabeling, the only configurations in which an agent whose
bundle contains only large items may fail to be envy-free, and hence
must use a large item as the EFX-removal item, are the following.
\begin{enumerate}[label=(\alph*)]
\item All items have type $(3,4)$ and
\[
 |R|=4q+3.
\]
Then agents $1,2$ receive $q$ and $q+1$ items, respectively, in either
order, while agents $3,4$ receive $q+1$ items each.  We are free to choose
which of agents $1,2$ receives $q+1$ items, and this agent is the unique agent who use a large item as the EFX-removal item.

\item The multiset $R$ consists of one item of type $(1,2)$ and
$4q+3$ items of type $(3,4)$.  One of agents $1,2$ receives the unique
type-$(1,2)$ item together with $q$ type-$(3,4)$ items, the other receives
$q+1$ type-$(3,4)$ items, and agents $3,4$ receive $q+1$ type-$(3,4)$
items each.  We are free to choose which of agents $1,2$ receives the unique
type-$(1,2)$ item.  The other agent is the unique agent who must use a large
item as the EFX-removal item.
\end{enumerate}

\item If every edge of $R$ has type $(1,2)$ or $(3,4)$, with
$x_{34}\ge x_{12}\ge 0$, then the allocation can be chosen so that agent $1$
is envy-free:
\[
 c_1(B_1)\le c_1(B_j)\qquad\forall j\in N.
\]
The same statement holds with agent $2$ in place of agent $1$.
\end{enumerate}
\end{lemma}

[Next verbatim source excerpt]

We will assume $r>2$ in the remaining part of this section.
For a chore $g$, let
\[
 S(g)=\{i\in N:c_i(g)=1\}
\]
be the set of agents for whom $g$ is small.  We partition the chores into
\[
M_{01}=\{g:|S(g)|\le 1\},\qquad
M_2=\{g:|S(g)|=2\},\qquad
M_{34}=\{g:|S(g)|\ge 3\}.
\]
For an item in $M_2$, we write $(i,j)$ for an item small exactly for agents
$i$ and $j$.  Thus the items in $M_2$ form a multigraph on the vertex set
$N$, where an edge $(i,j)$ is a chore small for precisely agents $i$ and $j$.

[Next verbatim source excerpt]

\subsection{The $M_2$ Multigraph Algorithm}
\label{subsec:2}
We now give a complete allocation procedure for the items in $M_2$.
Each item is small for some agents $i$ and $j$ and
large for the other two agents, and we say that the type of this item is $(i,j)$, or $ij$ for short.
In this section, we consider the case where the instance only has $M_{2}$ items.
We represent the instance by a
multigraph $R$ on $N=\{1,2,3,4\}$, where an item of type $(i,j)$ is an edge $(i,j)$.
We slightly abuse the notation by letting $R$ also represent the edge set, i.e., $R$ exactly corresponds to items in $M_2$.
Let $x_{ij}$ denote the multiplicity of edge $(i,j)$ and let $m_R=|R|$.

For $U\subseteq N$, define
\[
 \operatorname{inc}_R(U)=\left|\{e\in R:e\cap U\neq\emptyset\}\right|
\]
and
\[
 \Delta_R(U)=\frac{|U|m_R}{4}-\operatorname{inc}_R(U).
\]
We call $U$ \emph{deficient} if $\Delta_R(U)>0$.
\end{ReviewVerbatim}



\paragraph{Expanded PaperInterface specification}
\begin{ReviewVerbatim}
∀ (Item : Type) (r : ℝ) (cost : AppliedModelingLib.FairDivision.ChoreCost (Fin 4) Item) (chores : Finset Item),
  2 < r →
    AppliedModelingLib.FairDivision.IsOneOrRChoreCost cost r →
      (∀ item ∈ chores, AppliedModelingLib.FairDivision.IsSmallForExactlyTwo cost item) →
        ∃ (allocation : AppliedModelingLib.FairDivision.Allocation (Fin 4) Item),
          AppliedModelingLib.FairDivision.IsAllocationOf allocation chores ∧
            AppliedModelingLib.FairDivision.EFXForChores (AppliedModelingLib.FairDivision.additiveChoreCost cost)
                allocation ∧
              ((¬∃ (dominant : Finset (Fin 4)) (auxiliary : Finset (Fin 4)) (q : ℕ),
                      dominant.card = 2 ∧
                        auxiliary.card = 2 ∧
                          Disjoint dominant auxiliary ∧
                            ((∀ item ∈ chores, AppliedModelingLib.FairDivision.smallAgentSet cost item = dominant) ∧
                                chores.card = 4 * q + 3 ∨
                              ∃ exceptionalItem ∈ chores,
                                AppliedModelingLib.FairDivision.smallAgentSet cost exceptionalItem = auxiliary ∧
                                  (∀ item ∈ chores.erase exceptionalItem,
                                      AppliedModelingLib.FairDivision.smallAgentSet cost item = dominant) ∧
                                    (chores.erase exceptionalItem).card = 4 * q + 3)) →
                  (∀ (i : Fin 4),
                      (∃ item ∈ allocation i, AppliedModelingLib.FairDivision.IsSmallChore cost i item) →
                        ∀ (j : Fin 4),
                          AppliedModelingLib.FairDivision.additiveChoreCost cost i (allocation i) - 1 ≤
                            AppliedModelingLib.FairDivision.additiveChoreCost cost i (allocation j)) ∧
                    ∀ (i : Fin 4),
                      (∀ item ∈ allocation i, AppliedModelingLib.FairDivision.IsLargeChore cost r i item) →
                        ∀ (j : Fin 4),
                          AppliedModelingLib.FairDivision.additiveChoreCost cost i (allocation i) ≤
                            AppliedModelingLib.FairDivision.additiveChoreCost cost i (allocation j)) ∧
                (∀ (firstType secondType : Finset (Fin 4)),
                    firstType.card = 2 →
                      secondType.card = 2 →
                        Disjoint firstType secondType →
                          (∀ item ∈ chores,
                              AppliedModelingLib.FairDivision.smallAgentSet cost item = firstType ∨
                                AppliedModelingLib.FairDivision.smallAgentSet cost item = secondType) →
                            {item ∈ chores |
                                    AppliedModelingLib.FairDivision.smallAgentSet cost item = secondType}.card ≥
                                {item ∈ chores |
                                    AppliedModelingLib.FairDivision.smallAgentSet cost item = firstType}.card →
                              ∀ agent ∈ firstType,
                                ∃ (preferredAllocation : AppliedModelingLib.FairDivision.Allocation (Fin 4) Item),
                                  AppliedModelingLib.FairDivision.IsAllocationOf preferredAllocation chores ∧
                                    AppliedModelingLib.FairDivision.EFXForChores
                                        (AppliedModelingLib.FairDivision.additiveChoreCost cost) preferredAllocation ∧
                                      ∀ (other : Fin 4),
                                        AppliedModelingLib.FairDivision.additiveChoreCost cost agent
                                            (preferredAllocation agent) ≤
                                          AppliedModelingLib.FairDivision.additiveChoreCost cost agent
                                            (preferredAllocation other)) ∧
                  ∀ (dominant auxiliary : Finset (Fin 4)) (q : ℕ),
                    dominant.card = 2 →
                      auxiliary.card = 2 →
                        Disjoint dominant auxiliary →
                          ∀ special ∈ auxiliary,
                            ∀ companion ∈ auxiliary,
                              special ≠ companion →
                                ((∀ item ∈ chores, AppliedModelingLib.FairDivision.smallAgentSet cost item = dominant) ∧
                                      chores.card = 4 * q + 3 ∨
                                    ∃ exceptionalItem ∈ chores,
                                      AppliedModelingLib.FairDivision.smallAgentSet cost exceptionalItem = auxiliary ∧
                                        (∀ item ∈ chores.erase exceptionalItem,
                                            AppliedModelingLib.FairDivision.smallAgentSet cost item = dominant) ∧
                                          (chores.erase exceptionalItem).card = 4 * q + 3) →
                                  ∃ (exceptionalAllocation : AppliedModelingLib.FairDivision.Allocation (Fin 4) Item),
                                    AppliedModelingLib.FairDivision.IsAllocationOf exceptionalAllocation chores ∧
                                      AppliedModelingLib.FairDivision.EFXForChores
                                          (AppliedModelingLib.FairDivision.additiveChoreCost cost)
                                          exceptionalAllocation ∧
                                        (∀ (agent : Fin 4),
                                            agent ≠ special →
                                              ∀ (other : Fin 4),
                                                AppliedModelingLib.FairDivision.additiveChoreCost cost agent
                                                      (exceptionalAllocation agent) -
                                                    1 ≤
                                                  AppliedModelingLib.FairDivision.additiveChoreCost cost agent
                                                    (exceptionalAllocation other)) ∧
                                          (∀ item ∈ exceptionalAllocation special,
                                              AppliedModelingLib.FairDivision.IsLargeChore cost r special item) ∧
                                            (∀ (other : Fin 4),
                                                other ≠ companion →
                                                  AppliedModelingLib.FairDivision.additiveChoreCost cost special
                                                      (exceptionalAllocation special) ≤
                                                    AppliedModelingLib.FairDivision.additiveChoreCost cost special
                                                      (exceptionalAllocation other)) ∧
                                              AppliedModelingLib.FairDivision.additiveChoreCost cost special
                                                      (exceptionalAllocation special) -
                                                    r ≤
                                                  AppliedModelingLib.FairDivision.additiveChoreCost cost special
                                                    (exceptionalAllocation companion) ∧
                                                (((∀ item ∈ chores,
                                                      AppliedModelingLib.FairDivision.smallAgentSet cost item =
                                                        dominant) ∧
                                                    chores.card = 4 * q + 3 ∧
                                                      ∀ (agent : Fin 4),
                                                        Finset.card (exceptionalAllocation agent) =
                                                          if agent = companion then q else q + 1) ∨
                                                  ∃ exceptionalItem ∈ chores,
                                                    AppliedModelingLib.FairDivision.smallAgentSet cost exceptionalItem =
                                                        auxiliary ∧
                                                      (∀ item ∈ chores.erase exceptionalItem,
                                                          AppliedModelingLib.FairDivision.smallAgentSet cost item =
                                                            dominant) ∧
                                                        (chores.erase exceptionalItem).card = 4 * q + 3 ∧
                                                          exceptionalItem ∈ exceptionalAllocation companion ∧
                                                            ∀ (agent : Fin 4),
                                                              Finset.card (exceptionalAllocation agent) = q + 1)
\end{ReviewVerbatim}



\paragraph{Lean proof endpoint (built separately)}\begin{ReviewVerbatim}
HT26EFXChores.m2_efx_allocation_properties
\end{ReviewVerbatim}

\paragraph{Recorded source-to-Spec screening}
\textbf{Verdict:} \texttt{matches}
\\
{\raggedright
\textbf{Reason:} Encodes the ordinary-\allowbreak{}case small-\allowbreak{}removal/\allowbreak{}envy-\allowbreak{}free guarantees, both relabelled exceptional forms, selectable exceptional agent, and two-\allowbreak{}type preferred-\allowbreak{}agent clause.\allowbreak{}
\par}
\reviewmatch{Matches source input}
\noindent\textbf{Reviewer annotation}\par
\reviewerbox
\clearpage
\hypertarget{source-claim-c475426ec81f3892}{}
\section*{Review row 16}
\textbf{Source locator:} {\footnotesize\raggedright cited publication, lines 1709-\allowbreak{}1712\par}
\paragraph{Verbatim source input}
\begin{ReviewVerbatim}
\begin{lemma}\label{lem:easycombine}
    Suppose $P=(P_1,P_2,P_3,P_4)$ is a canonical allocation of $M_{01}$. Suppose some items of \(M_2\) are used for gap-filling, all assigned to short agents, such that \(P\) is extended to an envy-free allocation.
    If the residual of $M_2$ is of an exceptional configuration with exceptional agents $i$ and $j$ and both agents are long agents in the canonical allocation, then there exists an EFX allocation for $M_{01}\cup M_2$.
\end{lemma}

[Next verbatim source excerpt]

For each agent $i$, let
\[
 n_i=\left|\{g\in M_{01}:S(g)=\{i\}\}\right|
\]
be the number of $M_{01}$ items uniquely small for $i$.  Items in $M_{01}$
with $S(g)=\emptyset$ are all-large items.
Let $|M_{01}|=4a+b$ with $b\in\{0,1,2,3\}$.
As mentioned before, we will make sure each agent receives either $a$ items or $a+1$ items.

Fix quotas $q_i\in\{a,a+1\}$ with $\sum_i q_i=|M_{01}|$.
We consider allocations of $M_{01}$ with quotas $q=(q_1,q_2,q_3,q_4)$ in which each agent $i$ first receives as many of her uniquely-small items as possible, up to quota $q_i$, and the remaining slots are filled with the remaining items.
We will show that all such allocations satisfy EFX.
It is also important that we have the freedom to choose which agents receive $a$ items and which agents receive $a+1$ items.
This freedom enables us to append the allocation of $M_2$ after the allocation of $M_{01}$.

We begin by defining this type of allocations.


[Next verbatim source excerpt]

\subsection{The $M_2$ Multigraph Algorithm}
\label{subsec:2}
We now give a complete allocation procedure for the items in $M_2$.
Each item is small for some agents $i$ and $j$ and
large for the other two agents, and we say that the type of this item is $(i,j)$, or $ij$ for short.
In this section, we consider the case where the instance only has $M_{2}$ items.
We represent the instance by a
multigraph $R$ on $N=\{1,2,3,4\}$, where an item of type $(i,j)$ is an edge $(i,j)$.
We slightly abuse the notation by letting $R$ also represent the edge set, i.e., $R$ exactly corresponds to items in $M_2$.
Let $x_{ij}$ denote the multiplicity of edge $(i,j)$ and let $m_R=|R|$.

For $U\subseteq N$, define
\[
 \operatorname{inc}_R(U)=\left|\{e\in R:e\cap U\neq\emptyset\}\right|
\]
and
\[
 \Delta_R(U)=\frac{|U|m_R}{4}-\operatorname{inc}_R(U).
\]
We call $U$ \emph{deficient} if $\Delta_R(U)>0$.

[Next verbatim source excerpt]

\subsection{Concatenating the $M_{01}$ prefix and the $M_2$ allocation}
\label{subsec:combine}
We now finish the proof for $M_{01}\cup M_2$.  Write
\[
 |M_{01}|=4a+b,
 \qquad b\in\{0,1,2,3\}.
\]
We discuss the four values of $b$.
In each case and each of the sub-case, we will discuss how the concatenation is done and show that the resultant allocation is EFX.
We will exploit many features discussed earlier: the freedom to choose those agents that receive one more item when allocating $M_{01}$, the freedom to choose envy-free agents suggested in Lemma~\ref{lem:M2properties}, etc.

At a high level, we will try to use some items in $M_2$ for gap-filling such that the allocation of $M_{01}$ in the first phase becomes envy-free after gap-filling.
If the residual of $M_2$ is not of the two configurations in Lemma~\ref{lem:M2properties}(ii), we can allocate this residue so that every envy is up to a small item.
Lemma~\ref{lem:composition} then implies the combined allocation is EFX.
The tricky part is to try to make the residue of $M_2$ after gap-filling avoid the two configurations in Lemma~\ref{lem:M2properties}(ii).
However, there are hard cases where this is not always possible, and we need to handle these cases separately.

If, after the gap-filling, the residual $M_2$ items form one of the two configurations in Lemma~\ref{lem:M2properties}(ii), we will call this configuration \emph{exceptional}.
In an exceptional configuration, one edge $(i,j)$ has multiplicity at least $3$, and, other than those $(i,j)$-items, there can be at most one other item of type $(i',j')$ where $(i',j')$ is disjoint from $(i,j)$.
We will call the two agents $i',j'$ \emph{exceptional agents}.
By Lemma~\ref{lem:M2properties}(ii), for the two exceptional agents $i',j'$, one of them receives only large item in the residual of $M_2$ and envies the other by a large item.
We call this agent the \emph{disadvantageous agent}.
Lemma~\ref{lem:M2properties}(ii) tells us we are free to choose the disadvantageous agent from the two exceptional agents.

[Next verbatim source excerpt]

\begin{lemma}[Properties of the EFX allocations in Lemma~\ref{lem:M2}]\label{lem:M2properties}
The EFX allocation $B=(B_1,\ldots,B_4)$ ensured by Lemma~\ref{lem:M2} can be chosen to satisfy the
following additional properties.
\begin{enumerate}[label=(\roman*)]
\item Except in the two configurations described in~(ii), every possible envy is certified by removing a small item.  More precisely,
whenever $B_i$ contains an item that is small for agent $i$,
\[
 c_i(B_i)-1\le c_i(B_j)\qquad\forall j\in N,
\]
whereas every agent whose bundle contains only large items is envy-free:
\[
 c_i(B_i)\le c_i(B_j)\qquad\forall j\in N.
\]

\item Up to relabeling, the only configurations in which an agent whose
bundle contains only large items may fail to be envy-free, and hence
must use a large item as the EFX-removal item, are the following.
\begin{enumerate}[label=(\alph*)]
\item All items have type $(3,4)$ and
\[
 |R|=4q+3.
\]
Then agents $1,2$ receive $q$ and $q+1$ items, respectively, in either
order, while agents $3,4$ receive $q+1$ items each.  We are free to choose
which of agents $1,2$ receives $q+1$ items, and this agent is the unique agent who use a large item as the EFX-removal item.

\item The multiset $R$ consists of one item of type $(1,2)$ and
$4q+3$ items of type $(3,4)$.  One of agents $1,2$ receives the unique
type-$(1,2)$ item together with $q$ type-$(3,4)$ items, the other receives
$q+1$ type-$(3,4)$ items, and agents $3,4$ receive $q+1$ type-$(3,4)$
items each.  We are free to choose which of agents $1,2$ receives the unique
type-$(1,2)$ item.  The other agent is the unique agent who must use a large
item as the EFX-removal item.
\end{enumerate}

\item If every edge of $R$ has type $(1,2)$ or $(3,4)$, with

[Next verbatim source excerpt]


If, after the gap-filling, the residual $M_2$ items form one of the two configurations in Lemma~\ref{lem:M2properties}(ii), we will call this configuration \emph{exceptional}.
In an exceptional configuration, one edge $(i,j)$ has multiplicity at least $3$, and, other than those $(i,j)$-items, there can be at most one other item of type $(i',j')$ where $(i',j')$ is disjoint from $(i,j)$.
We will call the two agents $i',j'$ \emph{exceptional agents}.
By Lemma~\ref{lem:M2properties}(ii), for the two exceptional agents $i',j'$, one of them receives only large item in the residual of $M_2$ and envies the other by a large item.
We call this agent the \emph{disadvantageous agent}.
Lemma~\ref{lem:M2properties}(ii) tells us we are free to choose the disadvantageous agent from the two exceptional agents.
\end{ReviewVerbatim}



\paragraph{Expanded PaperInterface specification}
\begin{ReviewVerbatim}
∀ (Item : Type) (r : ℝ) (cost : AppliedModelingLib.FairDivision.ChoreCost (Fin 4) Item)
  (prefixChores m2Chores gapChores residueChores : Finset Item) (a : ℕ) (quota : Fin 4 → ℕ)
  (prefixAllocation gap : AppliedModelingLib.FairDivision.Allocation (Fin 4) Item),
  2 < r →
    AppliedModelingLib.FairDivision.IsOneOrRChoreCost cost r →
      Disjoint prefixChores m2Chores →
        (∀ item ∈ m2Chores, AppliedModelingLib.FairDivision.IsSmallForExactlyTwo cost item) →
          Disjoint gapChores residueChores →
            gapChores ∪ residueChores = m2Chores →
              AppliedModelingLib.FairDivision.IsCanonicalSmallChoreAllocation cost prefixChores quota prefixAllocation →
                (∀ item ∈ prefixChores, AppliedModelingLib.FairDivision.IsSmallForAtMostOne cost item) →
                  (∀ (agent : Fin 4), quota agent = a ∨ quota agent = a + 1) →
                    AppliedModelingLib.FairDivision.IsAllocationOf gap gapChores →
                      (∀ (agent : Fin 4), quota agent ≠ a → gap agent = ∅) →
                        (AppliedModelingLib.FairDivision.EnvyFreeForChores
                            (AppliedModelingLib.FairDivision.additiveChoreCost cost) fun (agent : Fin 4) =>
                            prefixAllocation agent ∪ gap agent) →
                          ∀ (exceptionalI exceptionalJ : Fin 4),
                            (∃ (dominant : Finset (Fin 4)) (auxiliary : Finset (Fin 4)) (q : ℕ),
                                dominant.card = 2 ∧
                                  auxiliary.card = 2 ∧
                                    Disjoint dominant auxiliary ∧
                                      exceptionalI ∈ auxiliary ∧
                                        exceptionalJ ∈ auxiliary ∧
                                          ((∀ item ∈ residueChores,
                                                AppliedModelingLib.FairDivision.smallAgentSet cost item = dominant) ∧
                                              residueChores.card = 4 * q + 3 ∨
                                            ∃ exceptionalItem ∈ residueChores,
                                              AppliedModelingLib.FairDivision.smallAgentSet cost exceptionalItem =
                                                  auxiliary ∧
                                                (∀ item ∈ residueChores.erase exceptionalItem,
                                                    AppliedModelingLib.FairDivision.smallAgentSet cost item =
                                                      dominant) ∧
                                                  (residueChores.erase exceptionalItem).card = 4 * q + 3)) →
                              exceptionalI ≠ exceptionalJ →
                                quota exceptionalI = a + 1 →
                                  quota exceptionalJ = a + 1 →
                                    ∃ (allocation : AppliedModelingLib.FairDivision.Allocation (Fin 4) Item),
                                      AppliedModelingLib.FairDivision.IsAllocationOf allocation
                                          (prefixChores ∪ m2Chores) ∧
                                        AppliedModelingLib.FairDivision.EFXForChores
                                          (AppliedModelingLib.FairDivision.additiveChoreCost cost) allocation
\end{ReviewVerbatim}



\paragraph{Lean proof endpoint (built separately)}\begin{ReviewVerbatim}
HT26EFXChores.exceptional_residue_combination
\end{ReviewVerbatim}

\paragraph{Recorded source-to-Spec screening}
\textbf{Verdict:} \texttt{matches}
\\
{\raggedright
\textbf{Reason:} Formalizes the canonical prefix, short-\allowbreak{}only envy-\allowbreak{}free gap filling, the two exceptional residual forms, two long exceptional agents, and EFX completion.\allowbreak{}
\par}
\reviewmatch{Matches source input}
\noindent\textbf{Reviewer annotation}\par
\reviewerbox
\clearpage
\hypertarget{source-claim-bf31a0d3fb856868}{}
\hypertarget{source-section-1743f0a28d68d6cc}{}
\section*{Appendix and supplement source claims}
\subsection*{Review row 17}
\textbf{Source locator:} {\footnotesize\raggedright cited publication, lines 2074-\allowbreak{}2076\par}
\paragraph{Verbatim source input}
\begin{ReviewVerbatim}
\begin{proposition}
    In an EFX allocation, no bundle contains two or more items from $A$.
\end{proposition}

[Next verbatim source excerpt]

In this section, we prove Theorem~\ref{thm:tri} for an arbitrary fixed $n\geq 4$.

\subsection{The construction}
Fix an integer $n\ge 4$. Partition the agents into two groups $T_1$ and $T_2$ such that
\[
    |T_1|=\left\lfloor \frac n2\right\rfloor,
    \qquad
    |T_2|=\left\lceil \frac n2\right\rceil.
\]
Let $t_1=|T_1|$ and $t_2=|T_2|$.
Let $s=2t_2+1$.
The instance contains $n-1+2s$ items that are partitioned into three groups $A,B,C$ where $A$ contains $n-1$ items and each of $B$ and $C$ contains $s$ items.

The three values in the tri-valued instances are given by
$$p=1,\quad q=s+2=2t_2+3,\quad\mbox{and}\quad r=\frac{s(q+1)}2=(2t_2+1)(t_2+2).$$


The costs are defined as follows. For every agent $i\in T_1$,
\[
    c_i(a)=r \quad(a\in A),
    \qquad
    c_i(b)=p \quad(b\in B),
    \qquad
    c_i(c)=q \quad(c\in C),
\]
and for every agent $i\in T_2$,
\[
    c_i(a)=r \quad(a\in A),
    \qquad
    c_i(b)=q \quad(b\in B),
    \qquad
    c_i(c)=p \quad(c\in C).
\]

For $n=4$, we have $t_1=t_2=2$, $s=5$, $q=7$, and $r=20$.
This resembles our construction in Sect.~\ref{sec:tri}.

\subsection{Non-existence of EFX allocations}
Let $X=(X_1,\ldots,X_n)$ be an arbitrary allocation. For a bundle $Y\subseteq A\cup B\cup C$, write
\[
    a(Y)=|Y\cap A|,
    \qquad
    b(Y)=|Y\cap B|,
    \qquad
    c(Y)=|Y\cap C|.
\]
All agents in $T_1$ have the same cost for a bundle $Y$, namely
\[
    P_1(Y)=r\cdot a(Y)+b(Y)+q\cdot c(Y),
\]
and all agents in $T_2$ have the same cost for $Y$, namely
\[
    P_2(Y)=r\cdot  a(Y)+q\cdot b(Y)+c(Y).
\]
Define
\[
    p_1=\min_{i\in \{1,\ldots,n\}} P_1(X_i),
    \qquad\mbox{and}\qquad
    p_2=\min_{i\in \{1,\ldots,n\}} P_2(X_i).
\]
\end{ReviewVerbatim}



\paragraph{Expanded PaperInterface specification}
\begin{ReviewVerbatim}
∀ (n : ℕ) (Item : Type) (r q : ℝ) (cost : AppliedModelingLib.FairDivision.ChoreCost (Fin n) Item)
  (chores A B C : Finset Item) (allocation : AppliedModelingLib.FairDivision.Allocation (Fin n) Item),
  4 ≤ n →
    q = ↑(2 * ((n + 1) / 2) + 1) + 2 →
      r = ↑(2 * ((n + 1) / 2) + 1) * (q + 1) / 2 →
        A.card = n - 1 →
          B.card = 2 * ((n + 1) / 2) + 1 →
            C.card = 2 * ((n + 1) / 2) + 1 →
              Disjoint A B →
                Disjoint A C →
                  Disjoint B C →
                    A ∪ B ∪ C = chores →
                      (∀ (agent : Fin n) (item : Item),
                          (item ∈ A → cost agent item = r) ∧
                            (item ∈ B → cost agent item = if ↑agent < n / 2 then 1 else q) ∧
                              (item ∈ C → cost agent item = if ↑agent < n / 2 then q else 1)) →
                        AppliedModelingLib.FairDivision.IsAllocationOf allocation chores →
                          AppliedModelingLib.FairDivision.EFXForChores
                              (AppliedModelingLib.FairDivision.additiveChoreCost cost) allocation →
                            ∀ (agent : Fin n), Finset.card (allocation agent ∩ A) ≤ 1
\end{ReviewVerbatim}



\paragraph{Lean proof endpoint (built separately)}\begin{ReviewVerbatim}
HT26EFXChores.appendix_no_two_a_items
\end{ReviewVerbatim}

\paragraph{Recorded source-to-Spec screening}
\textbf{Verdict:} \texttt{matches}
\\
{\raggedright
\textbf{Reason:} Exact Appendix construction, EFX allocation, and at-\allowbreak{}most-\allowbreak{}one-\allowbreak{}A-\allowbreak{}item conclusion.\allowbreak{}
\par}
\reviewmatch{Matches source input}
\noindent\textbf{Reviewer annotation}\par
\reviewerbox
\clearpage
\hypertarget{source-claim-18b6e05d84d21367}{}
\section*{Review row 18}
\textbf{Source locator:} {\footnotesize\raggedright cited publication, lines 2104-\allowbreak{}2106\par}
\paragraph{Verbatim source input}
\begin{ReviewVerbatim}
\begin{proposition}
    In an EFX allocation, $p_1\ge r$ and $p_2\ge r$.
\end{proposition}

[Next verbatim source excerpt]

In this section, we prove Theorem~\ref{thm:tri} for an arbitrary fixed $n\geq 4$.

\subsection{The construction}
Fix an integer $n\ge 4$. Partition the agents into two groups $T_1$ and $T_2$ such that
\[
    |T_1|=\left\lfloor \frac n2\right\rfloor,
    \qquad
    |T_2|=\left\lceil \frac n2\right\rceil.
\]
Let $t_1=|T_1|$ and $t_2=|T_2|$.
Let $s=2t_2+1$.
The instance contains $n-1+2s$ items that are partitioned into three groups $A,B,C$ where $A$ contains $n-1$ items and each of $B$ and $C$ contains $s$ items.

The three values in the tri-valued instances are given by
$$p=1,\quad q=s+2=2t_2+3,\quad\mbox{and}\quad r=\frac{s(q+1)}2=(2t_2+1)(t_2+2).$$


The costs are defined as follows. For every agent $i\in T_1$,
\[
    c_i(a)=r \quad(a\in A),
    \qquad
    c_i(b)=p \quad(b\in B),
    \qquad
    c_i(c)=q \quad(c\in C),
\]
and for every agent $i\in T_2$,
\[
    c_i(a)=r \quad(a\in A),
    \qquad
    c_i(b)=q \quad(b\in B),
    \qquad
    c_i(c)=p \quad(c\in C).
\]

For $n=4$, we have $t_1=t_2=2$, $s=5$, $q=7$, and $r=20$.
This resembles our construction in Sect.~\ref{sec:tri}.

\subsection{Non-existence of EFX allocations}
Let $X=(X_1,\ldots,X_n)$ be an arbitrary allocation. For a bundle $Y\subseteq A\cup B\cup C$, write
\[
    a(Y)=|Y\cap A|,
    \qquad
    b(Y)=|Y\cap B|,
    \qquad
    c(Y)=|Y\cap C|.
\]
All agents in $T_1$ have the same cost for a bundle $Y$, namely
\[
    P_1(Y)=r\cdot a(Y)+b(Y)+q\cdot c(Y),
\]
and all agents in $T_2$ have the same cost for $Y$, namely
\[
    P_2(Y)=r\cdot  a(Y)+q\cdot b(Y)+c(Y).
\]
Define
\[
    p_1=\min_{i\in \{1,\ldots,n\}} P_1(X_i),
    \qquad\mbox{and}\qquad
    p_2=\min_{i\in \{1,\ldots,n\}} P_2(X_i).
\]
\end{ReviewVerbatim}



\paragraph{Expanded PaperInterface specification}
\begin{ReviewVerbatim}
∀ (n : ℕ) (Item : Type) (r q : ℝ) (cost : AppliedModelingLib.FairDivision.ChoreCost (Fin n) Item)
  (chores A B C : Finset Item) (allocation : AppliedModelingLib.FairDivision.Allocation (Fin n) Item)
  (P1 P2 : AppliedModelingLib.FairDivision.Bundle Item → ℝ) (p1 p2 : ℝ),
  4 ≤ n →
    q = ↑(2 * ((n + 1) / 2) + 1) + 2 →
      r = ↑(2 * ((n + 1) / 2) + 1) * (q + 1) / 2 →
        A.card = n - 1 →
          B.card = 2 * ((n + 1) / 2) + 1 →
            C.card = 2 * ((n + 1) / 2) + 1 →
              Disjoint A B →
                Disjoint A C →
                  Disjoint B C →
                    A ∪ B ∪ C = chores →
                      (∀ (agent : Fin n) (item : Item),
                          (item ∈ A → cost agent item = r) ∧
                            (item ∈ B → cost agent item = if ↑agent < n / 2 then 1 else q) ∧
                              (item ∈ C → cost agent item = if ↑agent < n / 2 then q else 1)) →
                        AppliedModelingLib.FairDivision.IsAllocationOf allocation chores →
                          AppliedModelingLib.FairDivision.EFXForChores
                              (AppliedModelingLib.FairDivision.additiveChoreCost cost) allocation →
                            (∀ (bundle : AppliedModelingLib.FairDivision.Bundle Item),
                                P1 bundle =
                                  r * ↑(Finset.card (bundle ∩ A)) + ↑(Finset.card (bundle ∩ B)) +
                                    q * ↑(Finset.card (bundle ∩ C))) →
                              (∀ (bundle : AppliedModelingLib.FairDivision.Bundle Item),
                                  P2 bundle =
                                    r * ↑(Finset.card (bundle ∩ A)) + q * ↑(Finset.card (bundle ∩ B)) +
                                      ↑(Finset.card (bundle ∩ C))) →
                                (∀ (agent : Fin n),
                                    ↑agent < n / 2 →
                                      AppliedModelingLib.FairDivision.additiveChoreCost cost agent = P1) →
                                  (∀ (agent : Fin n),
                                      n / 2 ≤ ↑agent →
                                        AppliedModelingLib.FairDivision.additiveChoreCost cost agent = P2) →
                                    (∃ (agent : Fin n), P1 (allocation agent) = p1) →
                                      (∀ (agent : Fin n), p1 ≤ P1 (allocation agent)) →
                                        (∃ (agent : Fin n), P2 (allocation agent) = p2) →
                                          (∀ (agent : Fin n), p2 ≤ P2 (allocation agent)) → p1 ≥ r ∧ p2 ≥ r
\end{ReviewVerbatim}



\paragraph{Lean proof endpoint (built separately)}\begin{ReviewVerbatim}
HT26EFXChores.appendix_p1_p2_lower_bounds
\end{ReviewVerbatim}

\paragraph{Recorded source-to-Spec screening}
\textbf{Verdict:} \texttt{matches}
\\
{\raggedright
\textbf{Reason:} Encodes the stated P1/\allowbreak{}P2 definitions and concludes p1,p2 are each at least r.\allowbreak{}
\par}
\reviewmatch{Matches source input}
\noindent\textbf{Reviewer annotation}\par
\reviewerbox

\end{document}
